\documentclass[a4paper,11pt]{article}
        \usepackage[top=15mm,bottom=18mm,left=16mm,right=16mm]{geometry}
        \usepackage{fontspec}
        \usepackage{xeCJK}
        \setmainfont{Times New Roman}
        \setCJKmainfont{Yu Gothic}
        \setmonofont{Consolas}
        \setCJKmonofont{Yu Gothic}
        \usepackage{booktabs}
        \usepackage{longtable}
        \usepackage{tabularx}
        \usepackage{array}
        \usepackage{enumitem}
        \usepackage{hyperref}
        \usepackage{listings}
        \usepackage{xcolor}
        \usepackage{amsmath}
        \usepackage{amssymb}
        \hypersetup{
          colorlinks=true,
          linkcolor=blue,
          urlcolor=blue,
          pdftitle={車体物理・電気モデル本体},
          pdfauthor={Codex}
        }
        \lstset{
          basicstyle=\ttfamily\small,
          breaklines=true,
          columns=fullflexible,
          keepspaces=true,
          frame=single,
          framerule=0.2pt,
          rulecolor=\color{black},
          showstringspaces=false
        }
        \setlength{\parskip}{0.42em}
        \setlength{\parindent}{1em}
        \setlength{\emergencystretch}{3em}
        \renewcommand{\arraystretch}{1.15}
        \title{12. 車体物理・電気モデル本体}
        \author{solar\_ws0129-main}
        \date{2026-07-29}
        \begin{document}
        \maketitle
        \tableofcontents
        \clearpage

        \section{対象}
        \begin{tabularx}{\linewidth}{>{\raggedright\arraybackslash}p{0.18\linewidth} X}
        \toprule
        項目 & 内容 \\
        \midrule
        ファイル & \path|mpc_solarcar/model.py| \\
        ソースSHA-256 & \path|91a29cea04d184a73105566c514d824c8259be3a801df05ddfd82af7e8c863c1| \\
        種別 & Python \\
        区分 & model \\
        役割 & 空力、転がり、坂、PV、MPPT、drive/regen 効率、battery IV、SoC 更新を統合した vehicle model。 \\
        起動文脈 & planner と simulation の数理コア。 \\
        \bottomrule
        \end{tabularx}

        \section{このファイルを読む理由}
        空力、転がり、坂、PV、MPPT、drive/regen 効率、battery IV、SoC 更新を統合した vehicle model。

        \begin{itemize}[leftmargin=1.6em]
        \item SolarCarModel と Params が中心。
\item electrical\_balance が planner 側で最も多く呼ばれる。
\item resistive\_forces、battery\_iv、soc\_step が各所から再利用される。
        \end{itemize}

        \section{呼び出し関係}
        \subsection{どこから呼ばれるか}
        \begin{itemize}[leftmargin=1.6em]
        \item \path|mpc_node.py|
\item \path|solar_state_node.py|
\item \path|scripts/solar_sim.py|
\item \path|estimator.py|
        \end{itemize}

        \subsection{次に読むと理解が進むファイル}
        \begin{itemize}[leftmargin=1.6em]
        \item \path|mpc_solarcar/utils_maps.py|
        \end{itemize}

        \subsection{同一リポジトリ内の主要依存}
        \begin{itemize}[leftmargin=1.6em]
        \item \path|mpc_solarcar/utils_maps.py|
        \end{itemize}

        \section{ソース冒頭の把握}
        モジュール docstring または先頭意図の要約:

        モジュール docstring は明示されていない。

        \subsection{代表コード断片}
        \begin{lstlisting}[language=Python]
)

@dataclass
class Params:
    dt: float=600.0
    rho: float=1.18
    air_density_mode: str='constant'
    air_density_reference_pressure_pa: float=101325.0
    CdA: float=0.13
    Crr: float=0.002
    Crr_per_wheel: float=0.0
    m: float=250.0
    g: float=9.80665
    P_aux: float=60.0
    P_aux_stopped: float=60.0
    P_aux_night: float=0.0
    aux_night_ghi_threshold_wm2: float=20.0
    gear_eta: float=0.98
    gear_ratio: float=6.0
    wheel_radius: float=0.28
    wheel_count: int=4
    driven_wheel_count: int=2
\end{lstlisting}


        \section{主要構造の読み取り}
        主要クラスは Params, SolarCarModel。 主要関数は eff\_drive, eff\_regen, select\_drive\_mode, R\_int, pv\_balance, pv\_power\_mppt, charge\_efficiency, soc\_step。

        \subsection{主要クラス・関数}
            \begin{longtable}{p{0.07\linewidth} p{0.16\linewidth} p{0.25\linewidth} p{0.40\linewidth}}
            \toprule
            行 & 種別 & 名前 & 補足 \\
            \midrule
            17 & class & \path|Params| &  \\
73 & class & \path|SolarCarModel| &  \\
10 & function & \path|_is_symbolic| & args: x \\
74 & function & \path|__init__| & args: self, drive\_map\_path, regen\_map\_path, Rint\_map\_path, params, panel\_eff\_map\_path, mppt\_eff\_map\_path, drive\_map\_eco\_path, drive\_map\_power\_path, regen\_map\_eco\_path, regen\_map\_power\_path, ocv\_soc\_map\_path \\
124 & function & \path|eff_drive| & args: self, v\_ms, tau\_nm \\
137 & function & \path|eff_regen| & args: self, v\_ms, tau\_nm \\
150 & function & \path|_update_mode_limits| & args: self \\
158 & function & \path|_select_mode| & args: self, v\_ms, tau\_nm \\
169 & function & \path|select_drive_mode| & args: self, v\_ms, tau\_nm \\
172 & function & \path|R_int| & args: self, T\_C, z \\
179 & function & \path|pv_balance| & args: self, G\_poa, T\_cell\_C \\
222 & function & \path|pv_power_mppt| & args: self, G\_poa, T\_cell\_C \\
225 & function & \path|_scaled_slope_pct| & args: self, slope\_pct \\
228 & function & \path|charge_efficiency| & args: self, P\_pack \\
235 & function & \path|soc_step| & args: self, z, P\_pack, dt\_sec \\
262 & function & \path|ocv_from_soc| & args: self, z \\
271 & function & \path|load_ocv_map| & args: self, path \\
280 & function & \path|air_density| & args: self, ambient\_temp\_c, elevation\_m \\
297 & function & \path|resistive_forces| & args: self, v\_ms, slope\_pct, headwind\_ms \\
344 & function & \path|battery_iv| & args: self, P\_pack, z, Tbat\_C \\
368 & function & \path|mech_power| & args: self, v\_ms, slope\_pct, headwind\_ms, inertial\_power\_w \\
389 & function & \path|auxiliary_power| & args: self, aux\_power\_w \\
396 & function & \path|scheduled_auxiliary_power| & args: self \\
405 & function & \path|torque_from_mech| & args: self, P\_mech, v\_ms, wheel\_radius, ratio \\
416 & function & \path|electrical_balance| & args: self, v\_ms, slope\_pct, z, Tbat\_C, G\_poa, Tcell\_C, headwind\_ms, aux\_power\_w, inertial\_power\_w, ambient\_temp\_c, elevation\_m \\
            \bottomrule
            \end{longtable}

        \subsection{ファイルを上から読んだときの定義順}
            \begin{longtable}{p{0.07\linewidth} p{0.84\linewidth}}
            \toprule
            行 & 内容 \\
            \midrule
            3 & 例外処理を伴う try ブロックを実行する。 \\
10 & 関数 \_is\_symbolic を定義する。 \\
17 & クラス Params を定義する。 \\
73 & クラス SolarCarModel を定義する。 \\
            \bottomrule
            \end{longtable}

        \subsection{import 群の詳細}
            \begin{longtable}{p{0.07\linewidth} p{0.33\linewidth} p{0.50\linewidth}}
            \toprule
            行 & import 文 & このファイルで読む理由 \\
            \midrule
            1 & \path|import math| & Python標準のスカラー数値関数と定数を使うため。実際に参照する名前と行は使用位置欄で確認する。 このファイル内での主な使用位置は L329, L331, L336, L362。 \\
2 & \path|import numpy as np| & 数値配列、clip、補間、統計計算を行うため。 このファイル内での主な使用位置は L135, L148, L266, L269, L292, L347, L447。 \\
4 & \path|import casadi as ca| & casadi モジュールを利用するため。 このファイル内での主な使用位置は L6, L11, L12, L126, L128, L130, L139, L143, ...。 \\
7 & \path|from dataclasses import dataclass| & dataclasses から dataclass を読み込み、このファイルの処理を組み立てるため。 このファイル内での主な使用位置は L16。 \\
8 & \path|from .utils_maps import read_eff_map, read_Rint_map, read_map, bilinear_interp, read_1d_map| & 効率マップ・抵抗マップ読込補間 から read\_eff\_map, read\_Rint\_map, read\_map, bilinear\_interp, read\_1d\_map を読み込み、このファイルの内部処理を分担させるため。 実体ファイルは mpc\_solarcar/utils\_maps.py。 このファイル内での主な使用位置は L81, L82, L93, L95, L97, L99, L101, L106, ...。 \\
            \bottomrule
            \end{longtable}

        \section{このファイルを読む前に必要な基礎知識}
        次の章は、構文やROS用語を既知と仮定しないための説明である。

        \subsection{プログラム、プロセス、メモリ、オブジェクトを区別する}
ソースファイルはディスク上の文字列であり、それ自体は走っていない。Python実行可能プログラムがソースを読み、OSがその実行を一つのプロセスとして管理する。プロセスには仮想メモリ、開いているファイル、スレッド、終了コードなどが対応する。


\begin{lstlisting}
ソースファイル -> Pythonインタプリタが読み込む -> OS上のプロセス -> Pythonオブジェクトをメモリに生成 -> 関数やcallbackを実行
\end{lstlisting}


「メモリ上に生成する」とは、実行中プロセスが使う記憶領域に、その値の型、属性、参照関係を表す実体を用意することである。変数は実体そのものというより、そのオブジェクトを指す名前として理解するとPythonの代入が読みやすい。


\begin{lstlisting}[language=python]
node = MPCNode()
alias = node
# nodeとaliasは同じオブジェクトを参照する。
\end{lstlisting}


一つのプロセスには複数スレッドを持たせられる。スレッドは同じプロセスのメモリを共有するため、受信callbackと最適化callbackが同じself.zを書き換える場合は実行順と排他を検討する必要がある。

\paragraph{根拠資料}
\begin{itemize}[leftmargin=1.6em]
\item \href{https://docs.python.org/3/tutorial/classes.html}{Python公式チュートリアル: Classes}
\end{itemize}

\subsection{名前、代入、型、参照}
Pythonの代入文は、右辺を先に評価し、その結果のオブジェクトへ左辺の名前を結び付ける。`x = f()`では、まずfを呼び、その戻り値が得られてからxが更新される。


`float(...)`、`int(...)`、`bool(...)`は型名を呼び出して値を変換する。外部CSV、YAML、ROS parameterから来た値は型が期待どおりとは限らないため、このリポジトリでは明示変換が多い。


`None`は値が無いことを示す単一のオブジェクト、`math.nan`は浮動小数点値ではあるが有効な数値ではないことを示す。両者は用途が異なるため、`is None`と`math.isfinite`を使い分ける。

\paragraph{根拠資料}
\begin{itemize}[leftmargin=1.6em]
\item \href{https://docs.python.org/3/tutorial/classes.html}{Python公式チュートリアル: Classes}
\end{itemize}

\subsection{関数、引数、戻り値、スコープ}
`def`は関数本体をその場で実行する文ではない。関数オブジェクトを作り、その名前を現在の名前空間へ登録する。`f`は関数そのもの、`f()`はその関数を呼んだ結果である。


仮引数は関数定義側の受け取り名、実引数は呼出側から渡す値である。位置引数、キーワード引数、既定値付き引数、`*args`、`**kwargs`は渡し方が異なる。


\begin{lstlisting}[language=python]
def clip_speed(value, minimum=0.0, maximum=110.0):
    return min(maximum, max(minimum, value))

v = clip_speed(120.0, maximum=90.0)
\end{lstlisting}


関数内で作った通常の名前はローカルスコープに属する。メソッドが`self.z`を更新すればオブジェクトの状態が残るが、単に`z`を更新しただけならその呼出しのローカル値である。

\paragraph{根拠資料}
\begin{itemize}[leftmargin=1.6em]
\item \href{https://docs.python.org/3/tutorial/classes.html}{Python公式チュートリアル: Classes}
\end{itemize}

\subsection{module、package、import、console\_scripts}
Pythonファイルはmoduleとして読み込める。複数moduleをディレクトリにまとめたものがpackageである。`from .model import SolarCarModel`の先頭の点は、現在と同じpackage内のmodel moduleを指す相対importである。


import時にはファイルのトップレベル文が上から一度実行される。`def`や`class`は関数・クラスを登録するが、その本体の通常処理は呼び出すまで走らない。


setuptoolsの`console\_scripts`は、端末で使う実行可能名と`package.module:function`を対応付けるインストール時メタデータである。生成される実行用ラッパーはmoduleをimportし、指定関数を呼び、戻り値を終了コードとして扱う小さな入口である。


\begin{lstlisting}
ROS launchのexecutable名 -> install済みconsole script -> mpc_solarcar.mpc_nodeをimport -> main()を呼ぶ
\end{lstlisting}

\paragraph{根拠資料}
\begin{itemize}[leftmargin=1.6em]
\item \href{https://setuptools.pypa.io/en/latest/userguide/entry_point.html}{setuptools公式: Entry Points / Console Scripts}
\item \href{https://docs.python.org/3/tutorial/classes.html}{Python公式チュートリアル: Classes}
\end{itemize}

\subsection{例外、try/finally、with、資源解放}
例外は通常の戻り値とは別経路で異常を呼出元へ伝える。`try/except`は想定した異常を処理し、`finally`は成功・失敗にかかわらず後始末を行う。


`with open(...) as f:`はcontext managerを使い、ブロックを出るとファイルを閉じる。CSVやログの破損を避けるため、開いた資源の所有者と閉じる場所を明確にする。


`except Exception: pass`はノードを止めない利点がある一方、入力異常を隠して原因追跡を難しくする。安全に関係する値では、少なくとも頻度制限付き警告、異常カウンタ、fallback状態のpublishを検討する。



\subsection{class、独自クラス、インスタンス、self、継承}
クラスは、保持するデータと、そのデータを扱う関数を一つの型としてまとめる設計単位である。「独自クラス」とはPythonやROS 2が最初から用意した型ではなく、このプロジェクトが目的に合わせて定義した新しい型を指す。


`class MPCNode(Node)`は、rclpyのNodeを基底クラスとして継承し、Nodeが持つpublisher、subscription、timer、parameterなどの機能にMPC固有機能を追加する。丸括弧内は関数引数ではなく基底クラス指定である。


`MPCNode()`はクラスオブジェクトを呼び出して新しいインスタンスを作る。Pythonは新しい実体を用意し、その実体を第1引数selfとして`\_\_init\_\_`へ渡す。`self`は予約語ではなく慣習的な名前だが、この慣習を守ることでコードの意味が共有される。


\begin{lstlisting}[language=python]
class Counter:
    def __init__(self):
        self.value = 0

    def increment(self):
        self.value += 1

counter = Counter()
counter.increment()
\end{lstlisting}


`super().\_\_init\_\_(...)`は継承元の初期化処理を呼ぶ。MPCNodeではこれを省くとROSノードとして必要な内部実体が作られず、`create\_publisher`などを正しく使えない。

\paragraph{根拠資料}
\begin{itemize}[leftmargin=1.6em]
\item \href{https://docs.python.org/3/tutorial/classes.html}{Python公式チュートリアル: Classes}
\item \href{https://docs.ros.org/en/humble/Tutorials/Beginner-CLI-Tools/Understanding-ROS2-Nodes/Understanding-ROS2-Nodes.html}{ROS 2 Humble公式: Understanding nodes}
\end{itemize}

\subsection{先頭アンダースコア、\_\_init\_\_、名前付けの作法}
`self.\_step\_solar`の先頭1個のアンダースコアは「クラス内部で使う実装詳細」という慣習であり、アクセスを強制禁止する機構ではない。外部コードから呼べるが、公開APIとして依存しない意思表示である。


`\_\_init\_\_`の前後2個のアンダースコアはPythonが定めた特殊メソッド名である。クラス生成時に自動で呼ばれる。任意の名前に前後2個を付けて独自仕様を作ることは避ける。


`\_`で始まるローカル変数は、未使用または外部へ見せない意図を示す場合がある。例えば`\_base\_csv`は戻り値として受け取る必要はあるが、その後の処理では使わないことを読者へ示す。

\paragraph{根拠資料}
\begin{itemize}[leftmargin=1.6em]
\item \href{https://docs.python.org/3/tutorial/classes.html}{Python公式チュートリアル: Classes}
\end{itemize}

\subsection{デコレータと@dataclass}
`@名前`は、直後に定義した関数またはクラスを別の関数へ渡し、その結果で定義名を置き換えるデコレータ構文である。`@Class`という一般構文があるのではなく、`@dataclass`など実際のデコレータ名を書く。


`@dataclass`は型注釈付きフィールドから`\_\_init\_\_`、`\_\_repr\_\_`、比較処理などを自動生成する。物理パラメータや解析結果のように、名前付きデータを一まとまりで運ぶ用途に向く。


\begin{lstlisting}[language=python]
from dataclasses import dataclass

@dataclass
class State:
    soc: float
    temperature_c: float

state = State(soc=0.8, temperature_c=30.0)
\end{lstlisting}


自動生成は検証を自動で保証しない。単位、許容範囲、相互依存制約は`\_\_post\_init\_\_`や別の検証関数で確認する必要がある。

\paragraph{根拠資料}
\begin{itemize}[leftmargin=1.6em]
\item \href{https://docs.python.org/3/library/dataclasses.html}{Python公式ライブラリ: dataclasses}
\item \href{https://docs.python.org/3/tutorial/classes.html}{Python公式チュートリアル: Classes}
\end{itemize}

\subsection{list、dict、deque、NumPy配列、pandas DataFrame}
listは順序付き可変列、tupleは順序付きで通常変更しない列、dictはキーから値を引く対応表、dequeは両端追加・削除が効率的なキューである。どの構造を選ぶかはアクセス方法と更新方法で決まる。


NumPy配列は同種数値を連続的に扱い、要素ごとの演算、clip、補間、線形代数を簡潔に書く。shapeは各次元の要素数であり、速度系列なら通常`(N,)`、候補集団なら`(population, N)`となる。


pandas DataFrameは列名を持つ表である。CSV読込後は列型、欠損、単位、timezone、並び順、重複を明示的に処理しなければ、数値計算が動いても意味が誤る。



\subsection{車両力学、電力収支、効率map}
\[
F_{\mathrm{aero}}=\frac{1}{2}\rho C_dA(v+v_{\mathrm{wind}})^2,\quad F_{\mathrm{roll}}\approx mgC_{rr}\cos\theta,\quad F_{\mathrm{grade}}=mg\sin\theta
\]

車輪機械powerは概ね`P\_mech = F\_total * v + P\_inertia`で、駆動時と回生時で異なる効率mapを通してDC側powerへ変換する。mapは速度・torqueなどを軸にした実測または同定tableであり、範囲外の補間・clip規則もモデルの一部である。


\[
P_{\mathrm{pack}}=P_{\mathrm{drive,dc}}-P_{\mathrm{regen,dc}}+P_{\mathrm{aux}}-P_{\mathrm{pv}}
\]

符号規約は必ずコードで確認する。本プロジェクトでは正のpack powerを放電側としてSoCを減らす処理が中心であり、発電と回生はpack負荷を下げる方向に働く。



\subsection{SoC、内部抵抗、端子電圧、温度、MHE}
\[
V=V_{\mathrm{oc}}(z,T)-IR_{\mathrm{total}}(z,T,I),\qquad z_{k+1}=z_k-\frac{\eta(I,T)I\Delta t}{Q_{\mathrm{eff}}}
\]

SoCは直接完全には観測できないため、電流積算、OCV、端子電圧、温度、容量、効率を組み合わせる。内部抵抗はSoC、温度、電流方向で変わり、発熱と電圧制約の両方へ影響する。


MHEは有限時間窓の状態と観測誤差をまとめて最適化し、SoCや温度を推定する。古い測定や欠損値を同じ重みで使うと推定が壊れるため、timestamp freshnessと観測可用性を入口で確認する。

\paragraph{根拠資料}
\begin{itemize}[leftmargin=1.6em]
\item \href{https://docs.scipy.org/doc/scipy/reference/generated/scipy.optimize.minimize.html}{SciPy公式: scipy.optimize.minimize}
\end{itemize}

\subsection{freshness、filter、guard、fallback、fail-safe}
分散システムでは最後に受け取った値が現在も有効とは限らない。受信時刻とtimeoutからfreshnessを判定し、stale値を計画状態へ無条件に同期しない。


filterはnoiseと一時的な飛び値を抑えるが、遅れを生む。slew limitは指令変化率を制限する。安全guardはsolverのcost罰則とは別に、現在出力へ強制制約を適用する最後の防波堤である。


fallbackは失敗時の代替動作を事前に決める設計である。前回計画保持、物理に基づく決定論的入力、停止、低速制限などから、故障modeごとに選ぶ。fallback発生はstatusとlogへ残し、正常解と区別する。



\subsection{制御、状態、入力、モデル予測制御MPC}
制御対象の内部を表す状態をx、操作入力をu、外乱・予報をwとすると、離散モデルは`x[k+1] = f(x[k], u[k], w[k])`と書ける。ソーラーカーではSoC、電池温度、距離、速度などが状態候補、速度目標や駆動トルクが入力候補になる。


\[
\min_{u_0,\ldots,u_{N-1}} \sum_{k=0}^{N-1}\ell(x_k,u_k,w_k)+V_f(x_N)
\]

MPCは現在状態からNステップ先まで予測し、目的関数と制約を満たす入力系列を求める。ただし実際に適用するのは通常先頭入力だけで、次回は新しい実測状態から再び解く。これがreceding horizonである。


予測モデル、目的関数、制約、ホライズン、solver、初期値のどれかが変わると答えも変わる。「MPCを使う」だけでは仕様は決まらず、これらを単位付きで追う必要がある。

\paragraph{根拠資料}
\begin{itemize}[leftmargin=1.6em]
\item \href{https://docs.scipy.org/doc/scipy/reference/generated/scipy.optimize.minimize.html}{SciPy公式: scipy.optimize.minimize}
\end{itemize}

\subsection{CSV/YAMLのdata contractと検証}
data contractは列名、型、単位、timezone、欠損可否、並び順、重複、許容範囲、先頭行、encodingを事前に決めた仕様である。単にCSVとして読めることは、モデル入力として正しいことを意味しない。


同定用実測、map、route、forecast、stop、scheduleはそれぞれgrainが異なる。生成時にschema validation、物理範囲、時間単調性、route範囲、coverageを検査し、検査結果をartifactとして残す。


学習用データと独立検証データを分離し、RMSEだけでなくbias、時系列残差、energy積算誤差、終端SoC、温度・電圧制約、外挿領域を評価する。





        \section{関数・クラスを上から順に解説}
        \subsection{L10 関数 \_is\_symbolic}
            \begin{itemize}[leftmargin=1.6em]
              \item 行範囲: L10--L14
              \item 所属: module直下
              \item 定義: \path|_is_symbolic(x)|
              \item docstring: 明示されていない。
              \item このブロックが直接呼ぶ主な関数・メソッド: hasattr, is\_symbolic, isinstance
              \item 戻り値の要点: ca is not None and (isinstance(x, (ca.SX, ca.MX)) or (hasattr(x, 'is\_symbolic') and x.is\_symbolic()))
              \item 主なローカル代入名: 特になし。
              \item 読み取る主なインスタンス属性: 特になし。
              \item 更新する主なインスタンス属性: 特になし。
              \item 明示的に送出する例外: 明示的raiseはない。
              \item 制御構造の規模: 条件分岐 0、ループ 0、try 0
            \end{itemize}
            \paragraph{この定義を読むためのPython構文}
            \begin{itemize}[leftmargin=1.6em]
            \item 特別な構文上の補足は少ない。
            \end{itemize}
            \paragraph{上から順の処理}
            \begin{enumerate}[leftmargin=1.8em]
            \item ca is not None and (isinstance(x, (ca.SX, ca.MX)) or (hasattr(x, 'is\_symbolic') and x.is\_symbolic())) を返す。
            \end{enumerate}
            \paragraph{代表コード断片}
            \begin{lstlisting}[language=Python]
def _is_symbolic(x):
    return ca is not None and (
        isinstance(x, (ca.SX, ca.MX))
        or (hasattr(x, 'is_symbolic') and x.is_symbolic())
    )
\end{lstlisting}

\subsection{L17 クラス Params}
            \begin{itemize}[leftmargin=1.6em]
              \item 行範囲: L17--L71
              \item 所属: module直下
              \item 定義: \path|Params(bases=none)|
              \item docstring: 明示されていない。
              \item このブロックが直接呼ぶ主な関数・メソッド: 特に明示されていない。
              \item 戻り値の要点: 戻り値が無いか、return 文が明示されていない。
              \item 主なローカル代入名: 特になし。
              \item 読み取る主なインスタンス属性: 特になし。
              \item 更新する主なインスタンス属性: 特になし。
              \item 明示的に送出する例外: 明示的raiseはない。
              \item 制御構造の規模: 条件分岐 0、ループ 0、try 0
            \end{itemize}
            \paragraph{この定義を読むためのPython構文}
            \begin{itemize}[leftmargin=1.6em]
            \item class文は新しい型を定義する。丸括弧内は継承する基底クラスである。
\item @で始まる行は、定義したクラスを加工するdecoratorである。
            \end{itemize}
            \paragraph{上から順の処理}
            \begin{enumerate}[leftmargin=1.8em]
            \item dt に 600.0 を代入する。
\item rho に 1.18 を代入する。
\item air\_density\_mode に 'constant' を代入する。
\item air\_density\_reference\_pressure\_pa に 101325.0 を代入する。
\item CdA に 0.13 を代入する。
\item Crr に 0.002 を代入する。
\item Crr\_per\_wheel に 0.0 を代入する。
\item m に 250.0 を代入する。
\item g に 9.80665 を代入する。
\item P\_aux に 60.0 を代入する。
\item P\_aux\_stopped に 60.0 を代入する。
\item P\_aux\_night に 0.0 を代入する。
            \end{enumerate}
            \paragraph{代表コード断片}
            \begin{lstlisting}[language=Python]
class Params:
    dt: float=600.0
    rho: float=1.18
    air_density_mode: str='constant'
    air_density_reference_pressure_pa: float=101325.0
    CdA: float=0.13
    Crr: float=0.002
    Crr_per_wheel: float=0.0
    m: float=250.0
    g: float=9.80665
    P_aux: float=60.0
    P_aux_stopped: float=60.0
    P_aux_night: float=0.0
    aux_night_ghi_threshold_wm2: float=20.0
    gear_eta: float=0.98
    gear_ratio: float=6.0
    wheel_radius: float=0.28
    wheel_count: int=4
    driven_wheel_count: int=2
    motor_count: int=1
    motor_type: str='generic'
    inverter_eta: float=1.0
    pv_area: float=4.0
    pv_eta_ref: float=0.23
    pv_mu_p: float=-0.0045
    mppt_eta: float=0.95
    panel_gain: float=1.0
    # Calibration from the vehicle telemetry solar-power channel to actual
    # battery-side solar power. It is applied only at telemetry ingestion.
    solar_measurement_gain_to_pack: float=1.0
    # Aggregate battery-side limit of the installed MPPT channels. A value
    # at or below zero disables clipping for vehicles without declared data.
    pv_power_limit_w: float=0.0
    E_nom_Wh: float=3055.0
    Q_nom_Ah: float=0.0
...
\end{lstlisting}

\subsection{L73 クラス SolarCarModel}
            \begin{itemize}[leftmargin=1.6em]
              \item 行範囲: L73--L477
              \item 所属: module直下
              \item 定義: \path|SolarCarModel(bases=none)|
              \item docstring: 明示されていない。
              \item このブロックが直接呼ぶ主な関数・メソッド: 特に明示されていない。
              \item 戻り値の要点: 戻り値が無いか、return 文が明示されていない。
              \item 主なローカル代入名: 特になし。
              \item 読み取る主なインスタンス属性: 特になし。
              \item 更新する主なインスタンス属性: 特になし。
              \item 明示的に送出する例外: 明示的raiseはない。
              \item 制御構造の規模: 条件分岐 0、ループ 0、try 0
            \end{itemize}
            \paragraph{この定義を読むためのPython構文}
            \begin{itemize}[leftmargin=1.6em]
            \item class文は新しい型を定義する。丸括弧内は継承する基底クラスである。
\item try文は例外が起き得る処理と、異常時または終了時の経路を分ける。
            \end{itemize}
            \paragraph{上から順の処理}
            \begin{enumerate}[leftmargin=1.8em]
            \item 関数 \_\_init\_\_ を定義する。
\item 関数 eff\_drive を定義する。
\item 関数 eff\_regen を定義する。
\item 関数 \_update\_mode\_limits を定義する。
\item 関数 \_select\_mode を定義する。
\item 関数 select\_drive\_mode を定義する。
\item 関数 R\_int を定義する。
\item 関数 pv\_balance を定義する。
\item 関数 pv\_power\_mppt を定義する。
\item 関数 \_scaled\_slope\_pct を定義する。
\item 関数 charge\_efficiency を定義する。
\item 関数 soc\_step を定義する。
            \end{enumerate}
            \paragraph{代表コード断片}
            \begin{lstlisting}[language=Python]
class SolarCarModel:
    def __init__(self, drive_map_path, regen_map_path, Rint_map_path,
                 params=None, panel_eff_map_path=None, mppt_eff_map_path=None,
                 drive_map_eco_path=None, drive_map_power_path=None,
                 regen_map_eco_path=None, regen_map_power_path=None,
                 ocv_soc_map_path=None):
        self.p = params or Params()
        self.aux_power_override_w = None
        self.v_grid, self.tau_grid, self.Z_drv = read_eff_map(drive_map_path)
        self.v_gridR, self.tau_gridR, self.Z_reg = read_eff_map(regen_map_path)
        self.drive_mode = 'auto'
        self.drive_mode_default = 'eco'
        self.drive_mode_tau_margin = 0.0
        self.maps_drive = {
            'default': (self.v_grid, self.tau_grid, self.Z_drv),
        }
        self.maps_regen = {
            'default': (self.v_gridR, self.tau_gridR, self.Z_reg),
        }
        if drive_map_eco_path:
            self.maps_drive['eco'] = read_eff_map(drive_map_eco_path)
        if drive_map_power_path:
            self.maps_drive['power'] = read_eff_map(drive_map_power_path)
        if regen_map_eco_path:
            self.maps_regen['eco'] = read_eff_map(regen_map_eco_path)
        if regen_map_power_path:
            self.maps_regen['power'] = read_eff_map(regen_map_power_path)
        self._update_mode_limits()
        self.Tg, self.zg, self.Rmap = read_Rint_map(Rint_map_path)
        self.panel_eff_map = None
        self.mppt_eff_map = None
        if panel_eff_map_path:
            try:
                self.Gg, self.Tcg, self.Z_panel = read_map(panel_eff_map_path)
                self.panel_eff_map = True
...
\end{lstlisting}

\subsection{L74 関数 SolarCarModel.\_\_init\_\_}
            \begin{itemize}[leftmargin=1.6em]
              \item 行範囲: L74--L122
              \item 所属: \path|SolarCarModel|
              \item 定義: \path|__init__(self, drive_map_path, regen_map_path, Rint_map_path, params = None, panel_eff_map_path = None, mppt_eff_map_path = None, drive_map_eco_path = None, drive_map_power_path = None, regen_map_eco_path = None, regen_map_power_path = None, ocv_soc_map_path = None)|
              \item docstring: 明示されていない。
              \item このブロックが直接呼ぶ主な関数・メソッド: Params, \_update\_mode\_limits, read\_1d\_map, read\_Rint\_map, read\_eff\_map, read\_map
              \item 戻り値の要点: 戻り値が無いか、return 文が明示されていない。
              \item 主なローカル代入名: 特になし。
              \item 読み取る主なインスタンス属性: self.Z\_drv, self.Z\_reg, self.\_update\_mode\_limits, self.maps\_drive, self.maps\_regen, self.tau\_grid, self.tau\_gridR, self.v\_grid, self.v\_gridR
              \item 更新する主なインスタンス属性: self.Gg, self.Gm, self.Rmap, self.Tcg, self.Tg, self.Tm, self.Z\_drv, self.Z\_mppt, self.Z\_panel, self.Z\_reg, self.aux\_power\_override\_w, self.drive\_mode, self.drive\_mode\_default, self.drive\_mode\_tau\_margin, self.maps\_drive, self.maps\_regen, self.mppt\_eff\_map, self.ocv\_grid, self.ocv\_soc\_map, self.p, self.panel\_eff\_map, self.soc\_grid, self.tau\_grid, self.tau\_gridR
              \item 明示的に送出する例外: 明示的raiseはない。
              \item 制御構造の規模: 条件分岐 7、ループ 0、try 3
            \end{itemize}
            \paragraph{この定義を読むためのPython構文}
            \begin{itemize}[leftmargin=1.6em]
            \item 第1引数selfは、このメソッドを呼んだインスタンス自身を受け取る慣習名である。
\item try文は例外が起き得る処理と、異常時または終了時の経路を分ける。
            \end{itemize}
            \paragraph{上から順の処理}
            \begin{enumerate}[leftmargin=1.8em]
            \item self.p に params or Params() の結果を代入する。
\item self.aux\_power\_override\_w に None の結果を代入する。
\item (self.v\_grid, self.tau\_grid, self.Z\_drv) に read\_eff\_map(drive\_map\_path) の結果を代入する。
\item (self.v\_gridR, self.tau\_gridR, self.Z\_reg) に read\_eff\_map(regen\_map\_path) の結果を代入する。
\item self.drive\_mode に 'auto' の結果を代入する。
\item self.drive\_mode\_default に 'eco' の結果を代入する。
\item self.drive\_mode\_tau\_margin に 0.0 の結果を代入する。
\item self.maps\_drive に \{'default': (self.v\_grid, self.tau\_grid, self.Z\_drv)\} の結果を代入する。
\item self.maps\_regen に \{'default': (self.v\_gridR, self.tau\_gridR, self.Z\_reg)\} の結果を代入する。
\item 条件 drive\_map\_eco\_path を判定し、真なら内部処理を行う。
\item   self.maps\_drive['eco'] に read\_eff\_map(drive\_map\_eco\_path) の結果を代入する。
\item 条件 drive\_map\_power\_path を判定し、真なら内部処理を行う。
\item   self.maps\_drive['power'] に read\_eff\_map(drive\_map\_power\_path) の結果を代入する。
\item 条件 regen\_map\_eco\_path を判定し、真なら内部処理を行う。
\item   self.maps\_regen['eco'] に read\_eff\_map(regen\_map\_eco\_path) の結果を代入する。
\item 条件 regen\_map\_power\_path を判定し、真なら内部処理を行う。
\item   self.maps\_regen['power'] に read\_eff\_map(regen\_map\_power\_path) の結果を代入する。
\item self.\_update\_mode\_limits(...) を実行する。
\item (self.Tg, self.zg, self.Rmap) に read\_Rint\_map(Rint\_map\_path) の結果を代入する。
\item self.panel\_eff\_map に None の結果を代入する。
\item self.mppt\_eff\_map に None の結果を代入する。
\item 条件 panel\_eff\_map\_path を判定し、真なら内部処理を行う。
\item   例外処理を伴う try ブロックを実行する。
\item     (self.Gg, self.Tcg, self.Z\_panel) に read\_map(panel\_eff\_map\_path) の結果を代入する。
\item     self.panel\_eff\_map に True の結果を代入する。
\item     Exceptionを捕捉した場合:
\item     self.panel\_eff\_map に None の結果を代入する。
\item 条件 mppt\_eff\_map\_path を判定し、真なら内部処理を行う。
\item   例外処理を伴う try ブロックを実行する。
\item     (self.Gm, self.Tm, self.Z\_mppt) に read\_map(mppt\_eff\_map\_path) の結果を代入する。
\item     self.mppt\_eff\_map に True の結果を代入する。
\item     Exceptionを捕捉した場合:
\item     self.mppt\_eff\_map に None の結果を代入する。
\item self.ocv\_soc\_map に None の結果を代入する。
\item 条件 ocv\_soc\_map\_path を判定し、真なら内部処理を行う。
\item   例外処理を伴う try ブロックを実行する。
\item     (self.soc\_grid, self.ocv\_grid) に read\_1d\_map(ocv\_soc\_map\_path) の結果を代入する。
\item     self.ocv\_soc\_map に True の結果を代入する。
\item     Exceptionを捕捉した場合:
\item     self.ocv\_soc\_map に None の結果を代入する。
            \end{enumerate}
            \paragraph{代表コード断片}
            \begin{lstlisting}[language=Python]
    def __init__(self, drive_map_path, regen_map_path, Rint_map_path,
                 params=None, panel_eff_map_path=None, mppt_eff_map_path=None,
                 drive_map_eco_path=None, drive_map_power_path=None,
                 regen_map_eco_path=None, regen_map_power_path=None,
                 ocv_soc_map_path=None):
        self.p = params or Params()
        self.aux_power_override_w = None
        self.v_grid, self.tau_grid, self.Z_drv = read_eff_map(drive_map_path)
        self.v_gridR, self.tau_gridR, self.Z_reg = read_eff_map(regen_map_path)
        self.drive_mode = 'auto'
        self.drive_mode_default = 'eco'
        self.drive_mode_tau_margin = 0.0
        self.maps_drive = {
            'default': (self.v_grid, self.tau_grid, self.Z_drv),
        }
        self.maps_regen = {
            'default': (self.v_gridR, self.tau_gridR, self.Z_reg),
        }
        if drive_map_eco_path:
            self.maps_drive['eco'] = read_eff_map(drive_map_eco_path)
        if drive_map_power_path:
            self.maps_drive['power'] = read_eff_map(drive_map_power_path)
        if regen_map_eco_path:
            self.maps_regen['eco'] = read_eff_map(regen_map_eco_path)
        if regen_map_power_path:
            self.maps_regen['power'] = read_eff_map(regen_map_power_path)
        self._update_mode_limits()
        self.Tg, self.zg, self.Rmap = read_Rint_map(Rint_map_path)
        self.panel_eff_map = None
        self.mppt_eff_map = None
        if panel_eff_map_path:
            try:
                self.Gg, self.Tcg, self.Z_panel = read_map(panel_eff_map_path)
                self.panel_eff_map = True
            except Exception:
...
\end{lstlisting}

\subsection{L124 関数 SolarCarModel.eff\_drive}
            \begin{itemize}[leftmargin=1.6em]
              \item 行範囲: L124--L135
              \item 所属: \path|SolarCarModel|
              \item 定義: \path|eff_drive(self, v_ms, tau_nm)|
              \item docstring: 明示されていない。
              \item このブロックが直接呼ぶ主な関数・メソッド: \_is\_symbolic, \_select\_mode, abs, bilinear\_interp, clip, fabs, float, fmax, fmin, get, sqrt
              \item 戻り値の要点: float(np.clip(eff, 0.55, 0.99)) / ca.fmin(0.99, ca.fmax(0.55, eff))
              \item 主なローカル代入名: Z, eff, mode, t, tN, t\_grid, v, vN, v\_grid
              \item 読み取る主なインスタンス属性: self.\_select\_mode, self.maps\_drive, self.p
              \item 更新する主なインスタンス属性: 特になし。
              \item 明示的に送出する例外: 明示的raiseはない。
              \item 制御構造の規模: 条件分岐 1、ループ 0、try 0
            \end{itemize}
            \paragraph{この定義を読むためのPython構文}
            \begin{itemize}[leftmargin=1.6em]
            \item 第1引数selfは、このメソッドを呼んだインスタンス自身を受け取る慣習名である。
            \end{itemize}
            \paragraph{上から順の処理}
            \begin{enumerate}[leftmargin=1.8em]
            \item 条件 \_is\_symbolic(v\_ms) or \_is\_symbolic(tau\_nm) を判定し、真なら内部処理を行う。
\item   v に ca.fabs(v\_ms) の結果を代入する。
\item   t に ca.fabs(tau\_nm) の結果を代入する。
\item   vN に v / 35.0 の結果を代入する。
\item   tN に t / 60.0 の結果を代入する。
\item   eff に 0.92 - 0.08 * vN * vN - 0.06 * ca.sqrt(tN + 1e-09) の結果を代入する。
\item   eff に eff * float(self.p.drive\_eff\_scale) の結果を代入する。
\item   ca.fmin(0.99, ca.fmax(0.55, eff)) を返す。
\item mode に self.\_select\_mode(float(v\_ms), float(abs(tau\_nm))) の結果を代入する。
\item (v\_grid, t\_grid, Z) に self.maps\_drive.get(mode, self.maps\_drive['default']) の結果を代入する。
\item eff に float(bilinear\_interp(v\_grid, t\_grid, Z, float(v\_ms), float(abs(tau\_nm)))) の結果を代入する。
\item eff を Mult で更新する。
\item float(np.clip(eff, 0.55, 0.99)) を返す。
            \end{enumerate}
            \paragraph{代表コード断片}
            \begin{lstlisting}[language=Python]
    def eff_drive(self, v_ms, tau_nm):
        if _is_symbolic(v_ms) or _is_symbolic(tau_nm):
            v = ca.fabs(v_ms); t = ca.fabs(tau_nm)
            vN = v/35.0; tN = t/60.0
            eff = 0.92 - 0.08*vN*vN - 0.06*ca.sqrt(tN+1e-9)
            eff = eff * float(self.p.drive_eff_scale)
            return ca.fmin(0.99, ca.fmax(0.55, eff))
        mode = self._select_mode(float(v_ms), float(abs(tau_nm)))
        v_grid, t_grid, Z = self.maps_drive.get(mode, self.maps_drive['default'])
        eff = float(bilinear_interp(v_grid, t_grid, Z, float(v_ms), float(abs(tau_nm))))
        eff *= float(self.p.drive_eff_scale)
        return float(np.clip(eff, 0.55, 0.99))
\end{lstlisting}

\subsection{L137 関数 SolarCarModel.eff\_regen}
            \begin{itemize}[leftmargin=1.6em]
              \item 行範囲: L137--L148
              \item 所属: \path|SolarCarModel|
              \item 定義: \path|eff_regen(self, v_ms, tau_nm)|
              \item docstring: 明示されていない。
              \item このブロックが直接呼ぶ主な関数・メソッド: \_is\_symbolic, \_select\_mode, abs, bilinear\_interp, clip, fabs, float, fmax, fmin, get
              \item 戻り値の要点: float(np.clip(eff, 0.4, 0.95)) / ca.fmin(0.95, ca.fmax(0.4, eff))
              \item 主なローカル代入名: Z, eff, mode, t, tN, t\_grid, v, vN, v\_grid
              \item 読み取る主なインスタンス属性: self.\_select\_mode, self.maps\_regen, self.p
              \item 更新する主なインスタンス属性: 特になし。
              \item 明示的に送出する例外: 明示的raiseはない。
              \item 制御構造の規模: 条件分岐 1、ループ 0、try 0
            \end{itemize}
            \paragraph{この定義を読むためのPython構文}
            \begin{itemize}[leftmargin=1.6em]
            \item 第1引数selfは、このメソッドを呼んだインスタンス自身を受け取る慣習名である。
            \end{itemize}
            \paragraph{上から順の処理}
            \begin{enumerate}[leftmargin=1.8em]
            \item 条件 \_is\_symbolic(v\_ms) or \_is\_symbolic(tau\_nm) を判定し、真なら内部処理を行う。
\item   v に ca.fabs(v\_ms) の結果を代入する。
\item   t に ca.fabs(tau\_nm) の結果を代入する。
\item   vN に v / 35.0 の結果を代入する。
\item   tN に t / 60.0 の結果を代入する。
\item   eff に 0.7 + 0.12 * vN - 0.05 * (tN - 0.3) * (tN - 0.3) の結果を代入する。
\item   eff に eff * float(self.p.regen\_eff\_scale or self.p.drive\_eff\_scale) の結果を代入する。
\item   ca.fmin(0.95, ca.fmax(0.4, eff)) を返す。
\item mode に self.\_select\_mode(float(v\_ms), float(abs(tau\_nm))) の結果を代入する。
\item (v\_grid, t\_grid, Z) に self.maps\_regen.get(mode, self.maps\_regen['default']) の結果を代入する。
\item eff に float(bilinear\_interp(v\_grid, t\_grid, Z, float(v\_ms), float(abs(tau\_nm)))) の結果を代入する。
\item eff を Mult で更新する。
\item float(np.clip(eff, 0.4, 0.95)) を返す。
            \end{enumerate}
            \paragraph{代表コード断片}
            \begin{lstlisting}[language=Python]
    def eff_regen(self, v_ms, tau_nm):
        if _is_symbolic(v_ms) or _is_symbolic(tau_nm):
            v = ca.fabs(v_ms); t = ca.fabs(tau_nm)
            vN = v/35.0; tN = t/60.0
            eff = 0.70 + 0.12*vN - 0.05*(tN-0.3)*(tN-0.3)
            eff = eff * float(self.p.regen_eff_scale or self.p.drive_eff_scale)
            return ca.fmin(0.95, ca.fmax(0.40, eff))
        mode = self._select_mode(float(v_ms), float(abs(tau_nm)))
        v_grid, t_grid, Z = self.maps_regen.get(mode, self.maps_regen['default'])
        eff = float(bilinear_interp(v_grid, t_grid, Z, float(v_ms), float(abs(tau_nm))))
        eff *= float(self.p.regen_eff_scale or self.p.drive_eff_scale)
        return float(np.clip(eff, 0.40, 0.95))
\end{lstlisting}

\subsection{L150 関数 SolarCarModel.\_update\_mode\_limits}
            \begin{itemize}[leftmargin=1.6em]
              \item 行範囲: L150--L156
              \item 所属: \path|SolarCarModel|
              \item 定義: \path|_update_mode_limits(self)|
              \item docstring: 明示されていない。
              \item このブロックが直接呼ぶ主な関数・メソッド: float, items, max
              \item 戻り値の要点: 戻り値が無いか、return 文が明示されていない。
              \item 主なローカル代入名: \_, k, t\_grid
              \item 読み取る主なインスタンス属性: self.maps\_drive, self.tau\_max
              \item 更新する主なインスタンス属性: self.tau\_max
              \item 明示的に送出する例外: 明示的raiseはない。
              \item 制御構造の規模: 条件分岐 0、ループ 1、try 1
            \end{itemize}
            \paragraph{この定義を読むためのPython構文}
            \begin{itemize}[leftmargin=1.6em]
            \item 第1引数selfは、このメソッドを呼んだインスタンス自身を受け取る慣習名である。
\item try文は例外が起き得る処理と、異常時または終了時の経路を分ける。
            \end{itemize}
            \paragraph{上から順の処理}
            \begin{enumerate}[leftmargin=1.8em]
            \item self.tau\_max に \{\} の結果を代入する。
\item self.maps\_drive.items() を順に走査し、各要素を (k, (\_, t\_grid, \_)) に入れて処理する。
\item   例外処理を伴う try ブロックを実行する。
\item     self.tau\_max[k] に float(max(t\_grid)) の結果を代入する。
\item     Exceptionを捕捉した場合:
\item     self.tau\_max[k] に 0.0 の結果を代入する。
            \end{enumerate}
            \paragraph{代表コード断片}
            \begin{lstlisting}[language=Python]
    def _update_mode_limits(self):
        self.tau_max = {}
        for k, (_, t_grid, _) in self.maps_drive.items():
            try:
                self.tau_max[k] = float(max(t_grid))
            except Exception:
                self.tau_max[k] = 0.0
\end{lstlisting}

\subsection{L158 関数 SolarCarModel.\_select\_mode}
            \begin{itemize}[leftmargin=1.6em]
              \item 行範囲: L158--L167
              \item 所属: \path|SolarCarModel|
              \item 定義: \path|_select_mode(self, v_ms: float, tau_nm: float) -> str|
              \item docstring: 明示されていない。
              \item このブロックが直接呼ぶ主な関数・メソッド: float, get, lower, str
              \item 戻り値の要点: 'eco' if 'eco' in self.maps\_drive else 'default' / mode / 'power' if 'power' in self.maps\_drive else 'default'
              \item 主なローカル代入名: eco\_max, margin, mode
              \item 読み取る主なインスタンス属性: self.drive\_mode, self.drive\_mode\_tau\_margin, self.maps\_drive, self.tau\_max
              \item 更新する主なインスタンス属性: 特になし。
              \item 明示的に送出する例外: 明示的raiseはない。
              \item 制御構造の規模: 条件分岐 2、ループ 0、try 0
            \end{itemize}
            \paragraph{この定義を読むためのPython構文}
            \begin{itemize}[leftmargin=1.6em]
            \item 第1引数selfは、このメソッドを呼んだインスタンス自身を受け取る慣習名である。
\item `->`以後は戻り値の型注釈であり、実行時の値を自動変換する命令ではない。
            \end{itemize}
            \paragraph{上から順の処理}
            \begin{enumerate}[leftmargin=1.8em]
            \item mode に str(self.drive\_mode or 'default').lower() の結果を代入する。
\item 条件 mode in ('eco', 'power') を判定し、真なら内部処理を行う。
\item   mode を返す。
\item eco\_max に self.tau\_max.get('eco', self.tau\_max.get('default', 0.0)) の結果を代入する。
\item margin に float(self.drive\_mode\_tau\_margin or 0.0) の結果を代入する。
\item 条件 tau\_nm > eco\_max + margin を判定し、真なら内部処理を行う。
\item   'power' if 'power' in self.maps\_drive else 'default' を返す。
\item 'eco' if 'eco' in self.maps\_drive else 'default' を返す。
            \end{enumerate}
            \paragraph{代表コード断片}
            \begin{lstlisting}[language=Python]
    def _select_mode(self, v_ms: float, tau_nm: float) -> str:
        mode = str(self.drive_mode or 'default').lower()
        if mode in ('eco', 'power'):
            return mode
        # auto
        eco_max = self.tau_max.get('eco', self.tau_max.get('default', 0.0))
        margin = float(self.drive_mode_tau_margin or 0.0)
        if tau_nm > (eco_max + margin):
            return 'power' if 'power' in self.maps_drive else 'default'
        return 'eco' if 'eco' in self.maps_drive else 'default'
\end{lstlisting}

\subsection{L169 関数 SolarCarModel.select\_drive\_mode}
            \begin{itemize}[leftmargin=1.6em]
              \item 行範囲: L169--L170
              \item 所属: \path|SolarCarModel|
              \item 定義: \path|select_drive_mode(self, v_ms: float, tau_nm: float) -> str|
              \item docstring: 明示されていない。
              \item このブロックが直接呼ぶ主な関数・メソッド: \_select\_mode, abs
              \item 戻り値の要点: self.\_select\_mode(v\_ms, abs(tau\_nm))
              \item 主なローカル代入名: 特になし。
              \item 読み取る主なインスタンス属性: self.\_select\_mode
              \item 更新する主なインスタンス属性: 特になし。
              \item 明示的に送出する例外: 明示的raiseはない。
              \item 制御構造の規模: 条件分岐 0、ループ 0、try 0
            \end{itemize}
            \paragraph{この定義を読むためのPython構文}
            \begin{itemize}[leftmargin=1.6em]
            \item 第1引数selfは、このメソッドを呼んだインスタンス自身を受け取る慣習名である。
\item `->`以後は戻り値の型注釈であり、実行時の値を自動変換する命令ではない。
            \end{itemize}
            \paragraph{上から順の処理}
            \begin{enumerate}[leftmargin=1.8em]
            \item self.\_select\_mode(v\_ms, abs(tau\_nm)) を返す。
            \end{enumerate}
            \paragraph{代表コード断片}
            \begin{lstlisting}[language=Python]
    def select_drive_mode(self, v_ms: float, tau_nm: float) -> str:
        return self._select_mode(v_ms, abs(tau_nm))
\end{lstlisting}

\subsection{L172 関数 SolarCarModel.R\_int}
            \begin{itemize}[leftmargin=1.6em]
              \item 行範囲: L172--L177
              \item 所属: \path|SolarCarModel|
              \item 定義: \path|R_int(self, T_C, z)|
              \item docstring: 明示されていない。
              \item このブロックが直接呼ぶ主な関数・メソッド: \_is\_symbolic, bilinear\_interp, float
              \item 戻り値の要点: (R0 + R\_T + R\_z) * float(self.p.rint\_scale) / float(self.p.rint\_scale) * float(bilinear\_interp(self.Tg, self.zg, self.Rmap, T\_C, z))
              \item 主なローカル代入名: R0, R\_T, R\_z
              \item 読み取る主なインスタンス属性: self.Rmap, self.Tg, self.p, self.zg
              \item 更新する主なインスタンス属性: 特になし。
              \item 明示的に送出する例外: 明示的raiseはない。
              \item 制御構造の規模: 条件分岐 1、ループ 0、try 0
            \end{itemize}
            \paragraph{この定義を読むためのPython構文}
            \begin{itemize}[leftmargin=1.6em]
            \item 第1引数selfは、このメソッドを呼んだインスタンス自身を受け取る慣習名である。
            \end{itemize}
            \paragraph{上から順の処理}
            \begin{enumerate}[leftmargin=1.8em]
            \item 条件 \_is\_symbolic(T\_C) or \_is\_symbolic(z) を判定し、真なら内部処理を行う。
\item   R0 に 0.015 の結果を代入する。
\item   R\_T に 0.0002 * (25.0 - T\_C) の結果を代入する。
\item   R\_z に 0.01 * (1.0 - z) の結果を代入する。
\item   (R0 + R\_T + R\_z) * float(self.p.rint\_scale) を返す。
\item   上の条件が偽の場合:
\item   float(self.p.rint\_scale) * float(bilinear\_interp(self.Tg, self.zg, self.Rmap, T\_C, z)) を返す。
            \end{enumerate}
            \paragraph{代表コード断片}
            \begin{lstlisting}[language=Python]
    def R_int(self, T_C, z):
        if _is_symbolic(T_C) or _is_symbolic(z):
            R0=0.015; R_T=0.0002*(25.0-T_C); R_z=0.01*(1.0-z)
            return (R0+R_T+R_z) * float(self.p.rint_scale)
        else:
            return float(self.p.rint_scale) * float(bilinear_interp(self.Tg, self.zg, self.Rmap, T_C, z))
\end{lstlisting}

\subsection{L179 関数 SolarCarModel.pv\_balance}
            \begin{itemize}[leftmargin=1.6em]
              \item 行範囲: L179--L220
              \item 所属: \path|SolarCarModel|
              \item 定義: \path|pv_balance(self, G_poa, T_cell_C)|
              \item docstring: 明示されていない。
              \item このブロックが直接呼ぶ主な関数・メソッド: \_is\_symbolic, bilinear\_interp, dict, float, fmax, fmin, max, min
              \item 戻り値の要点: dict(eta\_panel=eta\_panel, eta\_mppt=eta\_mppt, P\_pv\_raw=P\_pv\_raw, P\_pv\_unlimited=P\_pv\_unlimited, P\_pv\_limit\_loss=P\_pv\_unlimited - P\_pv, P\_pv=P\_pv)
              \item 主なローカル代入名: P\_pv, P\_pv\_raw, P\_pv\_unlimited, eta\_mppt, eta\_panel, power\_limit\_w, symbolic
              \item 読み取る主なインスタンス属性: self.Gg, self.Gm, self.Tcg, self.Tm, self.Z\_mppt, self.Z\_panel, self.mppt\_eff\_map, self.p, self.panel\_eff\_map
              \item 更新する主なインスタンス属性: 特になし。
              \item 明示的に送出する例外: 明示的raiseはない。
              \item 制御構造の規模: 条件分岐 6、ループ 0、try 0
            \end{itemize}
            \paragraph{この定義を読むためのPython構文}
            \begin{itemize}[leftmargin=1.6em]
            \item 第1引数selfは、このメソッドを呼んだインスタンス自身を受け取る慣習名である。
            \end{itemize}
            \paragraph{上から順の処理}
            \begin{enumerate}[leftmargin=1.8em]
            \item symbolic に \_is\_symbolic(G\_poa) or \_is\_symbolic(T\_cell\_C) の結果を代入する。
\item 条件 symbolic を判定し、真なら内部処理を行う。
\item   eta\_panel に self.p.pv\_eta\_ref * (1.0 + self.p.pv\_mu\_p * (T\_cell\_C - 25.0)) の結果を代入する。
\item   eta\_panel に ca.fmax(0.0, eta\_panel) の結果を代入する。
\item   上の条件が偽の場合:
\item   条件 self.panel\_eff\_map を判定し、真なら内部処理を行う。
\item     eta\_panel に bilinear\_interp(self.Gg, self.Tcg, self.Z\_panel, float(G\_poa), float(T\_cell\_C)) の結果を代入する。
\item     eta\_panel に max(0.0, float(eta\_panel)) の結果を代入する。
\item     上の条件が偽の場合:
\item     eta\_panel に self.p.pv\_eta\_ref * (1.0 + self.p.pv\_mu\_p * (T\_cell\_C - 25.0)) の結果を代入する。
\item     eta\_panel に max(0.0, float(eta\_panel)) の結果を代入する。
\item eta\_panel を Mult で更新する。
\item P\_pv\_raw に eta\_panel * self.p.pv\_area * G\_poa の結果を代入する。
\item 条件 symbolic を判定し、真なら内部処理を行う。
\item   eta\_mppt に float(self.p.mppt\_eta) の結果を代入する。
\item   上の条件が偽の場合:
\item   条件 self.mppt\_eff\_map を判定し、真なら内部処理を行う。
\item     eta\_mppt に bilinear\_interp(self.Gm, self.Tm, self.Z\_mppt, float(G\_poa), float(T\_cell\_C)) の結果を代入する。
\item     eta\_mppt に max(0.0, float(eta\_mppt)) の結果を代入する。
\item     上の条件が偽の場合:
\item     eta\_mppt に float(self.p.mppt\_eta) の結果を代入する。
\item P\_pv\_unlimited に eta\_mppt * P\_pv\_raw の結果を代入する。
\item power\_limit\_w に max(0.0, float(self.p.pv\_power\_limit\_w)) の結果を代入する。
\item 条件 power\_limit\_w > 0.0 を判定し、真なら内部処理を行う。
\item   条件 symbolic を判定し、真なら内部処理を行う。
\item     P\_pv に ca.fmin(P\_pv\_unlimited, power\_limit\_w) の結果を代入する。
\item     上の条件が偽の場合:
\item     P\_pv に min(float(P\_pv\_unlimited), power\_limit\_w) の結果を代入する。
\item   上の条件が偽の場合:
\item   P\_pv に P\_pv\_unlimited の結果を代入する。
\item dict(eta\_panel=eta\_panel, eta\_mppt=eta\_mppt, P\_pv\_raw=P\_pv\_raw, P\_pv\_unlimited=P\_pv\_unlimited, P\_pv\_limit\_loss=P\_pv\_unlimited - P\_pv, P\_pv=P\_pv) を返す。
            \end{enumerate}
            \paragraph{代表コード断片}
            \begin{lstlisting}[language=Python]
    def pv_balance(self, G_poa, T_cell_C):
        symbolic = _is_symbolic(G_poa) or _is_symbolic(T_cell_C)
        if symbolic:
            # Map interpolation is numeric. Keep the symbolic optimization
            # path differentiable by using the map's calibrated reference
            # model, as is done for the symbolic drive-efficiency path.
            eta_panel = self.p.pv_eta_ref * (
                1.0 + self.p.pv_mu_p * (T_cell_C - 25.0)
            )
            eta_panel = ca.fmax(0.0, eta_panel)
        elif self.panel_eff_map:
            eta_panel = bilinear_interp(self.Gg, self.Tcg, self.Z_panel, float(G_poa), float(T_cell_C))
            eta_panel = max(0.0, float(eta_panel))
        else:
            eta_panel = self.p.pv_eta_ref*(1.0+self.p.pv_mu_p*(T_cell_C-25.0))
            eta_panel = max(0.0, float(eta_panel))
        eta_panel *= float(self.p.panel_gain)
        P_pv_raw = eta_panel*self.p.pv_area*G_poa
        if symbolic:
            eta_mppt = float(self.p.mppt_eta)
        elif self.mppt_eff_map:
            eta_mppt = bilinear_interp(self.Gm, self.Tm, self.Z_mppt, float(G_poa), float(T_cell_C))
            eta_mppt = max(0.0, float(eta_mppt))
        else:
            eta_mppt = float(self.p.mppt_eta)
        P_pv_unlimited = eta_mppt*P_pv_raw
        power_limit_w = max(0.0, float(self.p.pv_power_limit_w))
        if power_limit_w > 0.0:
            if symbolic:
                P_pv = ca.fmin(P_pv_unlimited, power_limit_w)
            else:
                P_pv = min(float(P_pv_unlimited), power_limit_w)
        else:
            P_pv = P_pv_unlimited
        return dict(
...
\end{lstlisting}

\subsection{L222 関数 SolarCarModel.pv\_power\_mppt}
            \begin{itemize}[leftmargin=1.6em]
              \item 行範囲: L222--L223
              \item 所属: \path|SolarCarModel|
              \item 定義: \path|pv_power_mppt(self, G_poa, T_cell_C)|
              \item docstring: 明示されていない。
              \item このブロックが直接呼ぶ主な関数・メソッド: pv\_balance
              \item 戻り値の要点: self.pv\_balance(G\_poa, T\_cell\_C)['P\_pv']
              \item 主なローカル代入名: 特になし。
              \item 読み取る主なインスタンス属性: self.pv\_balance
              \item 更新する主なインスタンス属性: 特になし。
              \item 明示的に送出する例外: 明示的raiseはない。
              \item 制御構造の規模: 条件分岐 0、ループ 0、try 0
            \end{itemize}
            \paragraph{この定義を読むためのPython構文}
            \begin{itemize}[leftmargin=1.6em]
            \item 第1引数selfは、このメソッドを呼んだインスタンス自身を受け取る慣習名である。
            \end{itemize}
            \paragraph{上から順の処理}
            \begin{enumerate}[leftmargin=1.8em]
            \item self.pv\_balance(G\_poa, T\_cell\_C)['P\_pv'] を返す。
            \end{enumerate}
            \paragraph{代表コード断片}
            \begin{lstlisting}[language=Python]
    def pv_power_mppt(self, G_poa, T_cell_C):
        return self.pv_balance(G_poa, T_cell_C)['P_pv']
\end{lstlisting}

\subsection{L225 関数 SolarCarModel.\_scaled\_slope\_pct}
            \begin{itemize}[leftmargin=1.6em]
              \item 行範囲: L225--L226
              \item 所属: \path|SolarCarModel|
              \item 定義: \path|_scaled_slope_pct(self, slope_pct)|
              \item docstring: 明示されていない。
              \item このブロックが直接呼ぶ主な関数・メソッド: float
              \item 戻り値の要点: slope\_pct * float(self.p.grade\_scale)
              \item 主なローカル代入名: 特になし。
              \item 読み取る主なインスタンス属性: self.p
              \item 更新する主なインスタンス属性: 特になし。
              \item 明示的に送出する例外: 明示的raiseはない。
              \item 制御構造の規模: 条件分岐 0、ループ 0、try 0
            \end{itemize}
            \paragraph{この定義を読むためのPython構文}
            \begin{itemize}[leftmargin=1.6em]
            \item 第1引数selfは、このメソッドを呼んだインスタンス自身を受け取る慣習名である。
            \end{itemize}
            \paragraph{上から順の処理}
            \begin{enumerate}[leftmargin=1.8em]
            \item slope\_pct * float(self.p.grade\_scale) を返す。
            \end{enumerate}
            \paragraph{代表コード断片}
            \begin{lstlisting}[language=Python]
    def _scaled_slope_pct(self, slope_pct):
        return slope_pct * float(self.p.grade_scale)
\end{lstlisting}

\subsection{L228 関数 SolarCarModel.charge\_efficiency}
            \begin{itemize}[leftmargin=1.6em]
              \item 行範囲: L228--L233
              \item 所属: \path|SolarCarModel|
              \item 定義: \path|charge_efficiency(self, P_pack) -> float|
              \item docstring: 明示されていない。
              \item このブロックが直接呼ぶ主な関数・メソッド: float
              \item 戻り値の要点: float(self.p.eta\_charge) if p\_pack < 0.0 else 1.0 / 1.0
              \item 主なローカル代入名: p\_pack
              \item 読み取る主なインスタンス属性: self.p
              \item 更新する主なインスタンス属性: 特になし。
              \item 明示的に送出する例外: 明示的raiseはない。
              \item 制御構造の規模: 条件分岐 0、ループ 0、try 1
            \end{itemize}
            \paragraph{この定義を読むためのPython構文}
            \begin{itemize}[leftmargin=1.6em]
            \item 第1引数selfは、このメソッドを呼んだインスタンス自身を受け取る慣習名である。
\item `->`以後は戻り値の型注釈であり、実行時の値を自動変換する命令ではない。
\item try文は例外が起き得る処理と、異常時または終了時の経路を分ける。
            \end{itemize}
            \paragraph{上から順の処理}
            \begin{enumerate}[leftmargin=1.8em]
            \item 例外処理を伴う try ブロックを実行する。
\item   p\_pack に float(P\_pack) の結果を代入する。
\item   Exceptionを捕捉した場合:
\item   1.0 を返す。
\item float(self.p.eta\_charge) if p\_pack < 0.0 else 1.0 を返す。
            \end{enumerate}
            \paragraph{代表コード断片}
            \begin{lstlisting}[language=Python]
    def charge_efficiency(self, P_pack) -> float:
        try:
            p_pack = float(P_pack)
        except Exception:
            return 1.0
        return float(self.p.eta_charge) if p_pack < 0.0 else 1.0
\end{lstlisting}

\subsection{L235 関数 SolarCarModel.soc\_step}
            \begin{itemize}[leftmargin=1.6em]
              \item 行範囲: L235--L260
              \item 所属: \path|SolarCarModel|
              \item 定義: \path!soc_step(self, z: float, P_pack: float, dt_sec: float, *, current_a: float | None = None, Tbat_C: float = 25.0) -> float!
              \item docstring: Advance SoC using the profile's declared state definition.

A positive ``Q\_nom\_Ah`` selects charge SoC and therefore coulomb
counting, which is consistent with an OCV-vs-SoC map.  Legacy
profiles that do not declare pack capacity retain energy integration
so existing non-solar profiles are not silently reinterpreted.
              \item このブロックが直接呼ぶ主な関数・メソッド: battery\_iv, charge\_efficiency, float, max
              \item 戻り値の要点: float(z) - eta * (float(P\_pack) * float(dt\_sec) / 3600.0) / max(float(self.p.E\_nom\_Wh), 1e-06) / float(z) - eta * current * float(dt\_sec) / (3600.0 * q\_nom\_ah)
              \item 主なローカル代入名: current, eta, q\_nom\_ah
              \item 読み取る主なインスタンス属性: self.battery\_iv, self.charge\_efficiency, self.p
              \item 更新する主なインスタンス属性: 特になし。
              \item 明示的に送出する例外: 明示的raiseはない。
              \item 制御構造の規模: 条件分岐 1、ループ 0、try 0
            \end{itemize}
            \paragraph{この定義を読むためのPython構文}
            \begin{itemize}[leftmargin=1.6em]
            \item 第1引数selfは、このメソッドを呼んだインスタンス自身を受け取る慣習名である。
\item `->`以後は戻り値の型注釈であり、実行時の値を自動変換する命令ではない。
            \end{itemize}
            \paragraph{上から順の処理}
            \begin{enumerate}[leftmargin=1.8em]
            \item eta に self.charge\_efficiency(P\_pack) の結果を代入する。
\item q\_nom\_ah に float(self.p.Q\_nom\_Ah) の結果を代入する。
\item 条件 q\_nom\_ah > 0.0 を判定し、真なら内部処理を行う。
\item   current に float(current\_a) if current\_a is not None else float(self.battery\_iv(P\_pack, z, Tbat\_C)['I']) の結果を代入する。
\item   float(z) - eta * current * float(dt\_sec) / (3600.0 * q\_nom\_ah) を返す。
\item float(z) - eta * (float(P\_pack) * float(dt\_sec) / 3600.0) / max(float(self.p.E\_nom\_Wh), 1e-06) を返す。
            \end{enumerate}
            \paragraph{代表コード断片}
            \begin{lstlisting}[language=Python]
    def soc_step(
        self,
        z: float,
        P_pack: float,
        dt_sec: float,
        *,
        current_a: float | None = None,
        Tbat_C: float = 25.0,
    ) -> float:
        """Advance SoC using the profile's declared state definition.

        A positive ``Q_nom_Ah`` selects charge SoC and therefore coulomb
        counting, which is consistent with an OCV-vs-SoC map.  Legacy
        profiles that do not declare pack capacity retain energy integration
        so existing non-solar profiles are not silently reinterpreted.
        """
        eta = self.charge_efficiency(P_pack)
        q_nom_ah = float(self.p.Q_nom_Ah)
        if q_nom_ah > 0.0:
            current = (
                float(current_a)
                if current_a is not None
                else float(self.battery_iv(P_pack, z, Tbat_C)["I"])
            )
            return float(z) - eta * current * float(dt_sec) / (3600.0 * q_nom_ah)
        return float(z) - eta * (float(P_pack) * float(dt_sec) / 3600.0) / max(float(self.p.E_nom_Wh), 1.0e-6)
\end{lstlisting}

\subsection{L262 関数 SolarCarModel.ocv\_from\_soc}
            \begin{itemize}[leftmargin=1.6em]
              \item 行範囲: L262--L269
              \item 所属: \path|SolarCarModel|
              \item 定義: \path|ocv_from_soc(self, z)|
              \item docstring: 明示されていない。
              \item このブロックが直接呼ぶ主な関数・メソッド: \_is\_symbolic, clip, float, fmax, fmin, interp
              \item 戻り値の要点: float(np.interp(zc, self.soc\_grid, self.ocv\_grid)) / self.p.V\_min + (self.p.V\_max - self.p.V\_min) * z\_clamped / self.p.V\_min + (self.p.V\_max - self.p.V\_min) * zc
              \item 主なローカル代入名: z\_clamped, zc
              \item 読み取る主なインスタンス属性: self.ocv\_grid, self.ocv\_soc\_map, self.p, self.soc\_grid
              \item 更新する主なインスタンス属性: 特になし。
              \item 明示的に送出する例外: 明示的raiseはない。
              \item 制御構造の規模: 条件分岐 2、ループ 0、try 0
            \end{itemize}
            \paragraph{この定義を読むためのPython構文}
            \begin{itemize}[leftmargin=1.6em]
            \item 第1引数selfは、このメソッドを呼んだインスタンス自身を受け取る慣習名である。
            \end{itemize}
            \paragraph{上から順の処理}
            \begin{enumerate}[leftmargin=1.8em]
            \item 条件 \_is\_symbolic(z) を判定し、真なら内部処理を行う。
\item   z\_clamped に ca.fmin(self.p.soc\_max, ca.fmax(self.p.soc\_min, z)) の結果を代入する。
\item   self.p.V\_min + (self.p.V\_max - self.p.V\_min) * z\_clamped を返す。
\item zc に float(np.clip(z, self.p.soc\_min, self.p.soc\_max)) の結果を代入する。
\item 条件 not self.ocv\_soc\_map を判定し、真なら内部処理を行う。
\item   self.p.V\_min + (self.p.V\_max - self.p.V\_min) * zc を返す。
\item float(np.interp(zc, self.soc\_grid, self.ocv\_grid)) を返す。
            \end{enumerate}
            \paragraph{代表コード断片}
            \begin{lstlisting}[language=Python]
    def ocv_from_soc(self, z):
        if _is_symbolic(z):
            z_clamped = ca.fmin(self.p.soc_max, ca.fmax(self.p.soc_min, z))
            return self.p.V_min + (self.p.V_max - self.p.V_min) * z_clamped
        zc = float(np.clip(z, self.p.soc_min, self.p.soc_max))
        if not self.ocv_soc_map:
            return self.p.V_min + (self.p.V_max - self.p.V_min) * zc
        return float(np.interp(zc, self.soc_grid, self.ocv_grid))
\end{lstlisting}

\subsection{L271 関数 SolarCarModel.load\_ocv\_map}
            \begin{itemize}[leftmargin=1.6em]
              \item 行範囲: L271--L278
              \item 所属: \path|SolarCarModel|
              \item 定義: \path|load_ocv_map(self, path: str) -> bool|
              \item docstring: 明示されていない。
              \item このブロックが直接呼ぶ主な関数・メソッド: read\_1d\_map
              \item 戻り値の要点: True / False
              \item 主なローカル代入名: 特になし。
              \item 読み取る主なインスタンス属性: 特になし。
              \item 更新する主なインスタンス属性: self.ocv\_grid, self.ocv\_soc\_map, self.soc\_grid
              \item 明示的に送出する例外: 明示的raiseはない。
              \item 制御構造の規模: 条件分岐 0、ループ 0、try 1
            \end{itemize}
            \paragraph{この定義を読むためのPython構文}
            \begin{itemize}[leftmargin=1.6em]
            \item 第1引数selfは、このメソッドを呼んだインスタンス自身を受け取る慣習名である。
\item `->`以後は戻り値の型注釈であり、実行時の値を自動変換する命令ではない。
\item try文は例外が起き得る処理と、異常時または終了時の経路を分ける。
            \end{itemize}
            \paragraph{上から順の処理}
            \begin{enumerate}[leftmargin=1.8em]
            \item 例外処理を伴う try ブロックを実行する。
\item   (self.soc\_grid, self.ocv\_grid) に read\_1d\_map(path) の結果を代入する。
\item   self.ocv\_soc\_map に True の結果を代入する。
\item   True を返す。
\item   Exceptionを捕捉した場合:
\item   self.ocv\_soc\_map に None の結果を代入する。
\item   False を返す。
            \end{enumerate}
            \paragraph{代表コード断片}
            \begin{lstlisting}[language=Python]
    def load_ocv_map(self, path: str) -> bool:
        try:
            self.soc_grid, self.ocv_grid = read_1d_map(path)
            self.ocv_soc_map = True
            return True
        except Exception:
            self.ocv_soc_map = None
            return False
\end{lstlisting}

\subsection{L280 関数 SolarCarModel.air\_density}
            \begin{itemize}[leftmargin=1.6em]
              \item 行範囲: L280--L295
              \item 所属: \path|SolarCarModel|
              \item 定義: \path|air_density(self, ambient_temp_c = None, elevation_m = 0.0)|
              \item docstring: Return dry-air density from the configured atmosphere model.

The ideal-gas-altitude mode uses the ISA tropospheric pressure law and
measured ambient temperature. ``rho`` remains the explicit fallback so
old profiles retain bit-for-bit constant-density behavior.
              \item このブロックが直接呼ぶ主な関数・メソッド: clip, float, lower, max, str
              \item 戻り値の要点: float(pressure\_pa / (287.05 * temperature\_k)) / float(self.p.rho) / float(self.p.rho)
              \item 主なローカル代入名: altitude, pressure\_pa, pressure\_ratio, temperature\_k
              \item 読み取る主なインスタンス属性: self.p
              \item 更新する主なインスタンス属性: 特になし。
              \item 明示的に送出する例外: 明示的raiseはない。
              \item 制御構造の規模: 条件分岐 2、ループ 0、try 0
            \end{itemize}
            \paragraph{この定義を読むためのPython構文}
            \begin{itemize}[leftmargin=1.6em]
            \item 第1引数selfは、このメソッドを呼んだインスタンス自身を受け取る慣習名である。
            \end{itemize}
            \paragraph{上から順の処理}
            \begin{enumerate}[leftmargin=1.8em]
            \item 条件 str(self.p.air\_density\_mode or 'constant').lower() != 'ideal\_gas\_altitude' を判定し、真なら内部処理を行う。
\item   float(self.p.rho) を返す。
\item 条件 ambient\_temp\_c is None を判定し、真なら内部処理を行う。
\item   float(self.p.rho) を返す。
\item temperature\_k に max(180.0, float(ambient\_temp\_c) + 273.15) の結果を代入する。
\item altitude に float(np.clip(elevation\_m, -500.0, 11000.0)) の結果を代入する。
\item pressure\_ratio に max(0.05, 1.0 - 0.0065 * altitude / 288.15) ** 5.255877 の結果を代入する。
\item pressure\_pa に max(1000.0, float(self.p.air\_density\_reference\_pressure\_pa)) * pressure\_ratio の結果を代入する。
\item float(pressure\_pa / (287.05 * temperature\_k)) を返す。
            \end{enumerate}
            \paragraph{代表コード断片}
            \begin{lstlisting}[language=Python]
    def air_density(self, ambient_temp_c=None, elevation_m=0.0):
        """Return dry-air density from the configured atmosphere model.

        The ideal-gas-altitude mode uses the ISA tropospheric pressure law and
        measured ambient temperature. ``rho`` remains the explicit fallback so
        old profiles retain bit-for-bit constant-density behavior.
        """
        if str(self.p.air_density_mode or 'constant').lower() != 'ideal_gas_altitude':
            return float(self.p.rho)
        if ambient_temp_c is None:
            return float(self.p.rho)
        temperature_k = max(180.0, float(ambient_temp_c) + 273.15)
        altitude = float(np.clip(elevation_m, -500.0, 11000.0))
        pressure_ratio = max(0.05, 1.0 - 0.0065 * altitude / 288.15) ** 5.255877
        pressure_pa = max(1000.0, float(self.p.air_density_reference_pressure_pa)) * pressure_ratio
        return float(pressure_pa / (287.05 * temperature_k))
\end{lstlisting}

\subsection{L297 関数 SolarCarModel.resistive\_forces}
            \begin{itemize}[leftmargin=1.6em]
              \item 行範囲: L297--L342
              \item 所属: \path|SolarCarModel|
              \item 定義: \path|resistive_forces(self, v_ms, slope_pct, headwind_ms = 0.0, *, ambient_temp_c = None, elevation_m = 0.0, air_density_kgm3 = None)|
              \item docstring: 明示されていない。
              \item このブロックが直接呼ぶ主な関数・メソッド: \_is\_symbolic, \_scaled\_slope\_pct, air\_density, atan, cos, dict, float, fmax, max, sin
              \item 戻り値の要点: dict(F\_aero=F\_aero, F\_roll=F\_roll, F\_grade=F\_grade, F\_total=F\_total, theta=theta, v\_rel\_ms=v\_rel, air\_density\_kgm3=float(rho), normal\_force\_n=N, Crr\_eff=Crr\_eff, slope\_scaled\_pct=float(self.\_scaled\_slope\_pct(slope\_pct))) / dict(F\_aero=F\_aero, F\_roll=F\_roll, F\_grade=F\_grade, F\_total=F\_total, theta=theta, v\_rel\_ms=v\_rel, air\_density\_kgm3=rho, normal\_force\_n=N, Crr\_eff=Crr\_eff, slope\_scaled\_pct=self.\_scaled\_slope\_pct(slope\_pct))
              \item 主なローカル代入名: Crr\_eff, F\_aero, F\_grade, F\_roll, F\_total, N, rho, theta, v\_rel
              \item 読み取る主なインスタンス属性: self.\_scaled\_slope\_pct, self.air\_density, self.p
              \item 更新する主なインスタンス属性: 特になし。
              \item 明示的に送出する例外: 明示的raiseはない。
              \item 制御構造の規模: 条件分岐 3、ループ 0、try 0
            \end{itemize}
            \paragraph{この定義を読むためのPython構文}
            \begin{itemize}[leftmargin=1.6em]
            \item 第1引数selfは、このメソッドを呼んだインスタンス自身を受け取る慣習名である。
            \end{itemize}
            \paragraph{上から順の処理}
            \begin{enumerate}[leftmargin=1.8em]
            \item rho に self.air\_density(ambient\_temp\_c, elevation\_m) if air\_density\_kgm3 is None else air\_density\_kgm3 の結果を代入する。
\item 条件 \_is\_symbolic(v\_ms) or \_is\_symbolic(slope\_pct) or \_is\_symbolic(headwind\_ms) を判定し、真なら内部処理を行う。
\item   v\_rel に ca.fmax(0.0, v\_ms + headwind\_ms) の結果を代入する。
\item   theta に ca.atan(self.\_scaled\_slope\_pct(slope\_pct) / 100.0) の結果を代入する。
\item   F\_aero に 0.5 * rho * self.p.CdA * v\_rel ** 2 の結果を代入する。
\item   N に self.p.m * self.p.g * ca.cos(theta) の結果を代入する。
\item   Crr\_eff に self.p.Crr の結果を代入する。
\item   条件 self.p.Crr\_per\_wheel and self.p.wheel\_count を判定し、真なら内部処理を行う。
\item     Crr\_eff に self.p.Crr\_per\_wheel * float(self.p.wheel\_count) の結果を代入する。
\item   F\_roll に Crr\_eff * N の結果を代入する。
\item   F\_grade に self.p.m * self.p.g * ca.sin(theta) の結果を代入する。
\item   F\_total に F\_aero + F\_roll + F\_grade の結果を代入する。
\item   dict(F\_aero=F\_aero, F\_roll=F\_roll, F\_grade=F\_grade, F\_total=F\_total, theta=theta, v\_rel\_ms=v\_rel, air\_density\_kgm3=rho, normal\_force\_n=N, Crr\_eff=Crr\_eff, slope\_scaled\_pct=self.\_scaled\_slope\_pct(slope\_pct)) を返す。
\item v\_rel に max(0.0, float(v\_ms) + float(headwind\_ms)) の結果を代入する。
\item theta に math.atan(float(self.\_scaled\_slope\_pct(slope\_pct)) / 100.0) の結果を代入する。
\item F\_aero に 0.5 * float(rho) * self.p.CdA * v\_rel ** 2 の結果を代入する。
\item N に self.p.m * self.p.g * math.cos(theta) の結果を代入する。
\item Crr\_eff に self.p.Crr の結果を代入する。
\item 条件 self.p.Crr\_per\_wheel and self.p.wheel\_count を判定し、真なら内部処理を行う。
\item   Crr\_eff に self.p.Crr\_per\_wheel * float(self.p.wheel\_count) の結果を代入する。
\item F\_roll に Crr\_eff * N の結果を代入する。
\item F\_grade に self.p.m * self.p.g * math.sin(theta) の結果を代入する。
\item F\_total に F\_aero + F\_roll + F\_grade の結果を代入する。
\item dict(F\_aero=F\_aero, F\_roll=F\_roll, F\_grade=F\_grade, F\_total=F\_total, theta=theta, v\_rel\_ms=v\_rel, air\_density\_kgm3=float(rho), normal\_force\_n=N, Crr\_eff=Crr\_eff, slope\_scaled\_pct=float(self.\_scaled\_slope\_pct(slope\_pct))) を返す。
            \end{enumerate}
            \paragraph{代表コード断片}
            \begin{lstlisting}[language=Python]
    def resistive_forces(
        self,
        v_ms,
        slope_pct,
        headwind_ms=0.0,
        *,
        ambient_temp_c=None,
        elevation_m=0.0,
        air_density_kgm3=None,
    ):
        rho = (
            self.air_density(ambient_temp_c, elevation_m)
            if air_density_kgm3 is None
            else air_density_kgm3
        )
        if _is_symbolic(v_ms) or _is_symbolic(slope_pct) or _is_symbolic(headwind_ms):
            v_rel = ca.fmax(0.0, v_ms + headwind_ms)
            theta = ca.atan(self._scaled_slope_pct(slope_pct) / 100.0)
            F_aero = 0.5 * rho * self.p.CdA * v_rel ** 2
            N = self.p.m * self.p.g * ca.cos(theta)
            Crr_eff = self.p.Crr
            if self.p.Crr_per_wheel and self.p.wheel_count:
                Crr_eff = self.p.Crr_per_wheel * float(self.p.wheel_count)
            F_roll = Crr_eff * N
            F_grade = self.p.m * self.p.g * ca.sin(theta)
            F_total = F_aero + F_roll + F_grade
            return dict(F_aero=F_aero, F_roll=F_roll, F_grade=F_grade,
                        F_total=F_total, theta=theta, v_rel_ms=v_rel,
                        air_density_kgm3=rho,
                        normal_force_n=N, Crr_eff=Crr_eff,
                        slope_scaled_pct=self._scaled_slope_pct(slope_pct))
        v_rel = max(0.0, float(v_ms) + float(headwind_ms))
        theta = math.atan(float(self._scaled_slope_pct(slope_pct)) / 100.0)
        F_aero = 0.5 * float(rho) * self.p.CdA * v_rel ** 2
        N = self.p.m * self.p.g * math.cos(theta)
...
\end{lstlisting}

\subsection{L344 関数 SolarCarModel.battery\_iv}
            \begin{itemize}[leftmargin=1.6em]
              \item 行範囲: L344--L366
              \item 所属: \path|SolarCarModel|
              \item 定義: \path|battery_iv(self, P_pack, z, Tbat_C)|
              \item docstring: 明示されていない。
              \item このブロックが直接呼ぶ主な関数・メソッド: R\_int, \_is\_symbolic, clip, dict, float, fmax, fmin, max, ocv\_from\_soc, sqrt
              \item 戻り値の要点: dict(I=I, V=V, OCV=OCV, Rint=Rint, Rline=Rline, Rpolarization=Rpolarization, Rtotal=Rtot, iv\_discriminant=disc)
              \item 主なローカル代入名: I, OCV, Rint, Rline, Rpolarization, Rtot, V, a, b, c, disc, symbolic, z\_rint
              \item 読み取る主なインスタンス属性: self.R\_int, self.ocv\_from\_soc, self.p
              \item 更新する主なインスタンス属性: 特になし。
              \item 明示的に送出する例外: 明示的raiseはない。
              \item 制御構造の規模: 条件分岐 1、ループ 0、try 0
            \end{itemize}
            \paragraph{この定義を読むためのPython構文}
            \begin{itemize}[leftmargin=1.6em]
            \item 第1引数selfは、このメソッドを呼んだインスタンス自身を受け取る慣習名である。
            \end{itemize}
            \paragraph{上から順の処理}
            \begin{enumerate}[leftmargin=1.8em]
            \item OCV に self.ocv\_from\_soc(z) の結果を代入する。
\item symbolic に \_is\_symbolic(P\_pack) or \_is\_symbolic(z) or \_is\_symbolic(Tbat\_C) の結果を代入する。
\item z\_rint に ca.fmin(0.95, ca.fmax(0.1, z)) if symbolic else float(np.clip(z, 0.1, 0.95)) の結果を代入する。
\item Rint に self.R\_int(Tbat\_C, z\_rint) の結果を代入する。
\item Rline に float(self.p.r\_line\_ohm) の結果を代入する。
\item Rpolarization に max(0.0, float(self.p.r\_polarization\_ohm)) の結果を代入する。
\item Rtot に Rint + Rline + Rpolarization の結果を代入する。
\item a に Rtot の結果を代入する。
\item b に -OCV の結果を代入する。
\item c に P\_pack の結果を代入する。
\item 条件 symbolic を判定し、真なら内部処理を行う。
\item   disc に ca.fmax(b * b - 4 * a * c, 0.0) の結果を代入する。
\item   I に (OCV - ca.sqrt(disc)) / (2 * Rtot) の結果を代入する。
\item   上の条件が偽の場合:
\item   disc に max(float(b * b - 4 * a * c), 0.0) の結果を代入する。
\item   I に (float(OCV) - math.sqrt(disc)) / (2 * float(Rtot)) の結果を代入する。
\item V に OCV - I * Rtot の結果を代入する。
\item dict(I=I, V=V, OCV=OCV, Rint=Rint, Rline=Rline, Rpolarization=Rpolarization, Rtotal=Rtot, iv\_discriminant=disc) を返す。
            \end{enumerate}
            \paragraph{代表コード断片}
            \begin{lstlisting}[language=Python]
    def battery_iv(self, P_pack, z, Tbat_C):
        OCV = self.ocv_from_soc(z)
        symbolic = _is_symbolic(P_pack) or _is_symbolic(z) or _is_symbolic(Tbat_C)
        z_rint = ca.fmin(0.95, ca.fmax(0.1, z)) if symbolic else float(np.clip(z, 0.1, 0.95))
        Rint = self.R_int(Tbat_C, z_rint)
        Rline = float(self.p.r_line_ohm)
        Rpolarization = max(0.0, float(self.p.r_polarization_ohm))
        # Planner intervals are much longer than the fitted polarization time
        # constant, so the stateless MPC model uses its steady-state resistance.
        Rtot = Rint + Rline + Rpolarization
        a = Rtot
        b = -OCV
        c = P_pack
        if symbolic:
            disc = ca.fmax(b * b - 4 * a * c, 0.0)
            I = (OCV - ca.sqrt(disc)) / (2 * Rtot)
        else:
            disc = max(float(b * b - 4 * a * c), 0.0)
            I = (float(OCV) - math.sqrt(disc)) / (2 * float(Rtot))
        V = OCV - I * Rtot
        return dict(I=I, V=V, OCV=OCV, Rint=Rint, Rline=Rline,
                    Rpolarization=Rpolarization,
                    Rtotal=Rtot, iv_discriminant=disc)
\end{lstlisting}

\subsection{L368 関数 SolarCarModel.mech\_power}
            \begin{itemize}[leftmargin=1.6em]
              \item 行範囲: L368--L387
              \item 所属: \path|SolarCarModel|
              \item 定義: \path|mech_power(self, v_ms, slope_pct, headwind_ms = 0.0, inertial_power_w = 0.0, *, ambient_temp_c = None, elevation_m = 0.0)|
              \item docstring: 明示されていない。
              \item このブロックが直接呼ぶ主な関数・メソッド: resistive\_forces
              \item 戻り値の要点: forces['F\_total'] * v\_ms + inertial\_power\_w
              \item 主なローカル代入名: forces
              \item 読み取る主なインスタンス属性: self.resistive\_forces
              \item 更新する主なインスタンス属性: 特になし。
              \item 明示的に送出する例外: 明示的raiseはない。
              \item 制御構造の規模: 条件分岐 0、ループ 0、try 0
            \end{itemize}
            \paragraph{この定義を読むためのPython構文}
            \begin{itemize}[leftmargin=1.6em]
            \item 第1引数selfは、このメソッドを呼んだインスタンス自身を受け取る慣習名である。
            \end{itemize}
            \paragraph{上から順の処理}
            \begin{enumerate}[leftmargin=1.8em]
            \item forces に self.resistive\_forces(v\_ms, slope\_pct, headwind\_ms, ambient\_temp\_c=ambient\_temp\_c, elevation\_m=elevation\_m) の結果を代入する。
\item forces['F\_total'] * v\_ms + inertial\_power\_w を返す。
            \end{enumerate}
            \paragraph{代表コード断片}
            \begin{lstlisting}[language=Python]
    def mech_power(
        self,
        v_ms,
        slope_pct,
        headwind_ms=0.0,
        inertial_power_w=0.0,
        *,
        ambient_temp_c=None,
        elevation_m=0.0,
    ):
        forces = self.resistive_forces(
            v_ms,
            slope_pct,
            headwind_ms,
            ambient_temp_c=ambient_temp_c,
            elevation_m=elevation_m,
        )
        # Aerodynamic force depends on air-relative speed, while wheel work is
        # force times ground speed. This matches the identification model.
        return forces['F_total'] * v_ms + inertial_power_w
\end{lstlisting}

\subsection{L389 関数 SolarCarModel.auxiliary\_power}
            \begin{itemize}[leftmargin=1.6em]
              \item 行範囲: L389--L394
              \item 所属: \path|SolarCarModel|
              \item 定義: \path|auxiliary_power(self, aux_power_w = None)|
              \item docstring: 明示されていない。
              \item このブロックが直接呼ぶ主な関数・メソッド: float, max
              \item 戻り値の要点: max(0.0, float(self.p.P\_aux)) / max(0.0, float(aux\_power\_w)) / max(0.0, float(self.aux\_power\_override\_w))
              \item 主なローカル代入名: 特になし。
              \item 読み取る主なインスタンス属性: self.aux\_power\_override\_w, self.p
              \item 更新する主なインスタンス属性: 特になし。
              \item 明示的に送出する例外: 明示的raiseはない。
              \item 制御構造の規模: 条件分岐 2、ループ 0、try 0
            \end{itemize}
            \paragraph{この定義を読むためのPython構文}
            \begin{itemize}[leftmargin=1.6em]
            \item 第1引数selfは、このメソッドを呼んだインスタンス自身を受け取る慣習名である。
            \end{itemize}
            \paragraph{上から順の処理}
            \begin{enumerate}[leftmargin=1.8em]
            \item 条件 aux\_power\_w is not None を判定し、真なら内部処理を行う。
\item   max(0.0, float(aux\_power\_w)) を返す。
\item 条件 self.aux\_power\_override\_w is not None を判定し、真なら内部処理を行う。
\item   max(0.0, float(self.aux\_power\_override\_w)) を返す。
\item max(0.0, float(self.p.P\_aux)) を返す。
            \end{enumerate}
            \paragraph{代表コード断片}
            \begin{lstlisting}[language=Python]
    def auxiliary_power(self, aux_power_w=None):
        if aux_power_w is not None:
            return max(0.0, float(aux_power_w))
        if self.aux_power_override_w is not None:
            return max(0.0, float(self.aux_power_override_w))
        return max(0.0, float(self.p.P_aux))
\end{lstlisting}

\subsection{L396 関数 SolarCarModel.scheduled\_auxiliary\_power}
            \begin{itemize}[leftmargin=1.6em]
              \item 行範囲: L396--L403
              \item 所属: \path|SolarCarModel|
              \item 定義: \path|scheduled_auxiliary_power(self, *, is_driving: bool, irradiance_wm2: float) -> float|
              \item docstring: 明示されていない。
              \item このブロックが直接呼ぶ主な関数・メソッド: auxiliary\_power, float, max
              \item 戻り値の要点: max(0.0, float(self.p.P\_aux\_stopped)) / self.auxiliary\_power() / max(0.0, float(self.p.P\_aux\_night)) / self.auxiliary\_power()
              \item 主なローカル代入名: 特になし。
              \item 読み取る主なインスタンス属性: self.aux\_power\_override\_w, self.auxiliary\_power, self.p
              \item 更新する主なインスタンス属性: 特になし。
              \item 明示的に送出する例外: 明示的raiseはない。
              \item 制御構造の規模: 条件分岐 3、ループ 0、try 0
            \end{itemize}
            \paragraph{この定義を読むためのPython構文}
            \begin{itemize}[leftmargin=1.6em]
            \item 第1引数selfは、このメソッドを呼んだインスタンス自身を受け取る慣習名である。
\item `->`以後は戻り値の型注釈であり、実行時の値を自動変換する命令ではない。
            \end{itemize}
            \paragraph{上から順の処理}
            \begin{enumerate}[leftmargin=1.8em]
            \item 条件 is\_driving を判定し、真なら内部処理を行う。
\item   self.auxiliary\_power() を返す。
\item 条件 float(irradiance\_wm2) <= max(0.0, float(self.p.aux\_night\_ghi\_threshold\_wm2)) を判定し、真なら内部処理を行う。
\item   max(0.0, float(self.p.P\_aux\_night)) を返す。
\item 条件 self.aux\_power\_override\_w is not None を判定し、真なら内部処理を行う。
\item   self.auxiliary\_power() を返す。
\item max(0.0, float(self.p.P\_aux\_stopped)) を返す。
            \end{enumerate}
            \paragraph{代表コード断片}
            \begin{lstlisting}[language=Python]
    def scheduled_auxiliary_power(self, *, is_driving: bool, irradiance_wm2: float) -> float:
        if is_driving:
            return self.auxiliary_power()
        if float(irradiance_wm2) <= max(0.0, float(self.p.aux_night_ghi_threshold_wm2)):
            return max(0.0, float(self.p.P_aux_night))
        if self.aux_power_override_w is not None:
            return self.auxiliary_power()
        return max(0.0, float(self.p.P_aux_stopped))
\end{lstlisting}

\subsection{L405 関数 SolarCarModel.torque\_from\_mech}
            \begin{itemize}[leftmargin=1.6em]
              \item 行範囲: L405--L414
              \item 所属: \path|SolarCarModel|
              \item 定義: \path|torque_from_mech(self, P_mech, v_ms, wheel_radius = None, ratio = None)|
              \item docstring: 明示されていない。
              \item このブロックが直接呼ぶ主な関数・メソッド: 特に明示されていない。
              \item 戻り値の要点: (T\_m, omega\_w * ratio)
              \item 主なローカル代入名: T\_m, T\_w, eps, omega\_w, ratio, wheel\_radius
              \item 読み取る主なインスタンス属性: self.p
              \item 更新する主なインスタンス属性: 特になし。
              \item 明示的に送出する例外: 明示的raiseはない。
              \item 制御構造の規模: 条件分岐 2、ループ 0、try 0
            \end{itemize}
            \paragraph{この定義を読むためのPython構文}
            \begin{itemize}[leftmargin=1.6em]
            \item 第1引数selfは、このメソッドを呼んだインスタンス自身を受け取る慣習名である。
            \end{itemize}
            \paragraph{上から順の処理}
            \begin{enumerate}[leftmargin=1.8em]
            \item 条件 wheel\_radius is None を判定し、真なら内部処理を行う。
\item   wheel\_radius に self.p.wheel\_radius の結果を代入する。
\item 条件 ratio is None を判定し、真なら内部処理を行う。
\item   ratio に self.p.gear\_ratio の結果を代入する。
\item eps に 0.001 の結果を代入する。
\item omega\_w に v\_ms / wheel\_radius の結果を代入する。
\item T\_w に P\_mech / (omega\_w + eps) の結果を代入する。
\item T\_m に T\_w / ratio の結果を代入する。
\item (T\_m, omega\_w * ratio) を返す。
            \end{enumerate}
            \paragraph{代表コード断片}
            \begin{lstlisting}[language=Python]
    def torque_from_mech(self, P_mech, v_ms, wheel_radius=None, ratio=None):
        if wheel_radius is None:
            wheel_radius = self.p.wheel_radius
        if ratio is None:
            ratio = self.p.gear_ratio
        eps=1e-3
        omega_w = v_ms/wheel_radius
        T_w = P_mech/(omega_w+eps)
        T_m = T_w/ratio
        return T_m, omega_w*ratio
\end{lstlisting}

\subsection{L416 関数 SolarCarModel.electrical\_balance}
            \begin{itemize}[leftmargin=1.6em]
              \item 行範囲: L416--L477
              \item 所属: \path|SolarCarModel|
              \item 定義: \path|electrical_balance(self, v_ms, slope_pct, z, Tbat_C, G_poa, Tcell_C, headwind_ms = 0.0, aux_power_w = None, inertial_power_w = 0.0, ambient_temp_c = None, elevation_m = 0.0)|
              \item docstring: 明示されていない。
              \item このブロックが直接呼ぶ主な関数・メソッド: \_is\_symbolic, auxiliary\_power, battery\_iv, charge\_efficiency, clip, dict, eff\_drive, eff\_regen, float, fmax, max, pv\_balance
              \item 戻り値の要点: dict(P\_pv=P\_pv, P\_pv\_raw=pv['P\_pv\_raw'], P\_pv\_unlimited=pv['P\_pv\_unlimited'], P\_pv\_limit\_loss=pv['P\_pv\_limit\_loss'], eta\_panel=pv['eta\_panel'], eta\_mppt=pv['eta\_mppt'], P\_mech=P\_mech\_wheel, P\_mech\_wheel=P\_mech\_wheel, P\_road\_load=P\_road\_load, P\_inertia=inertial\_power\_w, P\_aux=P\_aux, P\_pack=P\_pack, I=I, V=V, losses\_line=losses\_line, losses\_int=losses\_int, losses\_rint=losses\_rint, losses\_polarization=losses\_polarization, OCV=OCV, Rint=Rint, Rline=Rline, Rpolarization=Rpolarization, Rtotal=Rtot, iv\_discriminant=disc, eta\_charge=eta\_charge, P\_dc\_to\_drv=P\_dc\_to\_drv, P\_reg\_to\_dc=P\_reg\_to\_dc, eff\_drv=eff\_drv, eff\_reg=eff\_reg, torque\_drive\_nm=Tm\_drv, torque\_regen\_nm=Tm\_reg, omega\_motor\_radps=omega\_m, omega\_wheel\_radps=omega\_m / max(float(self.p.gear\_ratio), 1e-09))
              \item 主なローカル代入名: I, OCV, P\_aux, P\_dc\_to\_drv, P\_mech\_neg, P\_mech\_pos, P\_mech\_wheel, P\_pack, P\_pv, P\_reg\_to\_dc, P\_road\_load, Rint, Rline, Rpolarization, Rtot, Tm\_drv, Tm\_reg, V, \_, disc
              \item 読み取る主なインスタンス属性: self.auxiliary\_power, self.battery\_iv, self.charge\_efficiency, self.eff\_drive, self.eff\_regen, self.p, self.pv\_balance, self.resistive\_forces, self.torque\_from\_mech
              \item 更新する主なインスタンス属性: 特になし。
              \item 明示的に送出する例外: 明示的raiseはない。
              \item 制御構造の規模: 条件分岐 1、ループ 0、try 0
            \end{itemize}
            \paragraph{この定義を読むためのPython構文}
            \begin{itemize}[leftmargin=1.6em]
            \item 第1引数selfは、このメソッドを呼んだインスタンス自身を受け取る慣習名である。
            \end{itemize}
            \paragraph{上から順の処理}
            \begin{enumerate}[leftmargin=1.8em]
            \item pv に self.pv\_balance(G\_poa, Tcell\_C) の結果を代入する。
\item P\_pv に pv['P\_pv'] の結果を代入する。
\item road\_forces に self.resistive\_forces(v\_ms, slope\_pct, headwind\_ms, ambient\_temp\_c=ambient\_temp\_c, elevation\_m=elevation\_m) の結果を代入する。
\item P\_road\_load に road\_forces['F\_total'] * v\_ms の結果を代入する。
\item P\_mech\_wheel に P\_road\_load + inertial\_power\_w の結果を代入する。
\item symbolic に \_is\_symbolic(P\_mech\_wheel) or \_is\_symbolic(z) or \_is\_symbolic(Tbat\_C) or \_is\_symbolic(inertial\_power\_w) の結果を代入する。
\item 条件 symbolic を判定し、真なら内部処理を行う。
\item   P\_mech\_pos に ca.fmax(P\_mech\_wheel, 0.0) の結果を代入する。
\item   P\_mech\_neg に ca.fmax(-P\_mech\_wheel, 0.0) の結果を代入する。
\item   上の条件が偽の場合:
\item   P\_mech\_pos に max(float(P\_mech\_wheel), 0.0) の結果を代入する。
\item   P\_mech\_neg に max(-float(P\_mech\_wheel), 0.0) の結果を代入する。
\item (Tm\_drv, omega\_m) に self.torque\_from\_mech(P\_mech\_pos, v\_ms) の結果を代入する。
\item eff\_drv に self.eff\_drive(v\_ms, Tm\_drv) の結果を代入する。
\item P\_dc\_to\_drv に P\_mech\_pos / (eff\_drv * self.p.gear\_eta * self.p.inverter\_eta) の結果を代入する。
\item (Tm\_reg, \_) に self.torque\_from\_mech(P\_mech\_neg, v\_ms) の結果を代入する。
\item eff\_reg に self.eff\_regen(v\_ms, Tm\_reg) の結果を代入する。
\item regen\_utilization に float(np.clip(self.p.regen\_utilization, 0.0, 1.0)) の結果を代入する。
\item P\_reg\_to\_dc に regen\_utilization * eff\_reg * self.p.gear\_eta * self.p.inverter\_eta * P\_mech\_neg の結果を代入する。
\item P\_aux に self.auxiliary\_power(aux\_power\_w) の結果を代入する。
\item P\_pack に P\_dc\_to\_drv - P\_reg\_to\_dc + P\_aux - P\_pv の結果を代入する。
\item iv に self.battery\_iv(P\_pack, z, Tbat\_C) の結果を代入する。
\item OCV に iv['OCV'] の結果を代入する。
\item Rint に iv['Rint'] の結果を代入する。
\item Rline に iv['Rline'] の結果を代入する。
\item Rpolarization に iv['Rpolarization'] の結果を代入する。
\item Rtot に iv['Rtotal'] の結果を代入する。
\item disc に iv['iv\_discriminant'] の結果を代入する。
\item I に iv['I'] の結果を代入する。
\item V に iv['V'] の結果を代入する。
\item losses\_line に I * I * Rline の結果を代入する。
\item losses\_rint に I * I * Rint の結果を代入する。
\item losses\_polarization に I * I * Rpolarization の結果を代入する。
\item losses\_int に losses\_rint + losses\_polarization の結果を代入する。
\item eta\_charge に self.charge\_efficiency(P\_pack) の結果を代入する。
\item dict(P\_pv=P\_pv, P\_pv\_raw=pv['P\_pv\_raw'], P\_pv\_unlimited=pv['P\_pv\_unlimited'], P\_pv\_limit\_loss=pv['P\_pv\_limit\_loss'], eta\_panel=pv['eta\_panel'], eta\_mppt=pv['eta\_mppt'], P\_mech=P\_mech\_wheel, P\_mech\_wheel=P\_mech\_wheel, P\_road\_load=P\_road\_load, P\_inertia=inertial\_power\_w, P\_aux=P\_aux, P\_pack=P\_pack, I=I, V=V, losses\_line=losses\_line, losses\_int=losses\_int, losses\_rint=losses\_rint, losses\_polarization=losses\_polarization, OCV=OCV, Rint=Rint, Rline=Rline, Rpolarization=Rpolarization, Rtotal=Rtot, iv\_discriminant=disc, eta\_charge=eta\_charge, P\_dc\_to\_drv=P\_dc\_to\_drv, P\_reg\_to\_dc=P\_reg\_to\_dc, eff\_drv=eff\_drv, eff\_reg=eff\_reg, torque\_drive\_nm=Tm\_drv, torque\_regen\_nm=Tm\_reg, omega\_motor\_radps=omega\_m, omega\_wheel\_radps=omega\_m / max(float(self.p.gear\_ratio), 1e-09)) を返す。
            \end{enumerate}
            \paragraph{代表コード断片}
            \begin{lstlisting}[language=Python]
    def electrical_balance(self, v_ms, slope_pct, z, Tbat_C, G_poa, Tcell_C,
                           headwind_ms=0.0, aux_power_w=None, inertial_power_w=0.0,
                           ambient_temp_c=None, elevation_m=0.0):
        pv = self.pv_balance(G_poa, Tcell_C)
        P_pv = pv['P_pv']
        road_forces = self.resistive_forces(
            v_ms,
            slope_pct,
            headwind_ms,
            ambient_temp_c=ambient_temp_c,
            elevation_m=elevation_m,
        )
        P_road_load = road_forces['F_total'] * v_ms
        P_mech_wheel = P_road_load + inertial_power_w
        symbolic = (
            _is_symbolic(P_mech_wheel)
            or _is_symbolic(z)
            or _is_symbolic(Tbat_C)
            or _is_symbolic(inertial_power_w)
        )
        if symbolic:
            P_mech_pos = ca.fmax(P_mech_wheel, 0.0)
            P_mech_neg = ca.fmax(-P_mech_wheel, 0.0)
        else:
            P_mech_pos = max(float(P_mech_wheel), 0.0)
            P_mech_neg = max(-float(P_mech_wheel), 0.0)
        Tm_drv, omega_m = self.torque_from_mech(P_mech_pos, v_ms)
        eff_drv = self.eff_drive(v_ms, Tm_drv)
        P_dc_to_drv = P_mech_pos/(eff_drv*self.p.gear_eta*self.p.inverter_eta)
        Tm_reg, _ = self.torque_from_mech(P_mech_neg, v_ms)
        eff_reg = self.eff_regen(v_ms, Tm_reg)
        regen_utilization = float(np.clip(self.p.regen_utilization, 0.0, 1.0))
        P_reg_to_dc = regen_utilization*eff_reg*self.p.gear_eta*self.p.inverter_eta*P_mech_neg
        P_aux = self.auxiliary_power(aux_power_w)
        P_pack = P_dc_to_drv - P_reg_to_dc + P_aux - P_pv
...
\end{lstlisting}











        \section{処理の流れ}
        \begin{enumerate}[leftmargin=1.7em]
        \item ソース中の主要関数を通じて処理を進める。
        \end{enumerate}

        \section{このファイルの位置づけ}
        この資料群では、\path|mpc_solarcar/model.py| を
        \textbf{model}の中核として扱う。
        したがって、ここで扱う処理は単独で完結せず、
        前段の profile / launch / telemetry と後段の planner / logger / report へ繋がる。

        \end{document}
